%-----------------------------------------------------------------------------------
%   PACKAGES
%-----------------------------------------------------------------------------------

\documentclass[a4paper,11pt,abstract=true]{scrartcl}

\usepackage[margin=.9in]{geometry}


\usepackage{mlmodern}


\usepackage{amsmath,amsthm,amssymb}
\usepackage{mathtools,mathrsfs}
\usepackage{centernot}
\usepackage{wrapfig}

\usepackage[dvipsnames]{xcolor} % moved up: must be loaded before tikz and tcolorbox

\usepackage{tikz,tikz-cd,asymptote}
\usetikzlibrary{positioning, fit, calc}

\usepackage[most]{tcolorbox}
\usepackage{enumerate}
\usepackage[a]{esvect}
\usepackage{hyperref} % loaded last among the packages to avoid clashes


\makeatletter
\def\@seccntformat#1{\textcolor{Red}{\hspace*{0.2em}$\textbf{\S}$\hspace{-0.3em} \csname the#1\endcsname}\hspace{0.5em}}
\makeatother

\usepackage[headsepline]{scrlayer-scrpage}
\clearpairofpagestyles{}
\ihead{Vector Spaces}
\ohead{\pagemark}
\setkomafont{pageheadfoot}{\normalfont}
\setkomafont{pagenumber}{\normalfont}
\pagestyle{scrheadings}

\setlength{\parindent}{0cm}

%-----------------------------------------------------------------------------------
%   COMMANDS
%-----------------------------------------------------------------------------------

\DeclarePairedDelimiter{\floor}{\lfloor}{\rfloor}

\newcommand{\nonadj}{\,\neg\textsf{adj}\,}
\newcommand{\adj}{\,\textsf{adj}\,}
\newcommand{\RR}{\mathbb{R}}
\newcommand{\ZZ}{\mathbb{Z}}
\newcommand{\NN}{\mathbb{N}}
\newcommand{\CC}{\mathbb{C}}
\newcommand{\QQ}{\mathbb{Q}}
\newcommand{\sgn}{\text{sgn\hspace{0.5mm}}}
\newcommand{\osc}{\text{osc\hspace{0.5mm}}}
\newcommand{\foralls}{\hspace{2mm}\forall}
\newcommand{\rimply}{\Rightarrow}
\newcommand{\limply}{\Leftarrow}
\newcommand{\underdash}{\rule{0.25cm}{0.1mm}\hspace{0.5mm}}
\newcommand{\undersmalldash}{\rule{0.18cm}{0.1mm}\hspace{0.3mm}}
\newcommand{\s}{\square}
\newcommand{\xr}{\xrightarrow}

\newcommand{\genstirlingI}[3]{%
  \genfrac{\lbrack}{\rbrack}{0pt}{#1}{#2}{#3}%
}
\newcommand{\genstirlingII}[3]{%
  \genfrac{\{}{\}}{0pt}{#1}{#2}{#3}%
}
\newcommand{\genlah}[3]{%
  \genfrac{\lfloor}{\rfloor}{0pt}{#1}{#2}{#3}%
}
\newcommand{\stirlingI}[2]{\genstirlingI{}{#1}{#2}}
\newcommand{\stirlingII}[2]{\genstirlingII{}{#1}{#2}}
\newcommand{\lah}[2]{\genlah{}{#1}{#2}}

%-----------------------------------------------------------------------------------
%   ENVIRONMENTS
%-----------------------------------------------------------------------------------

\newtheoremstyle{dotlessP}{}{}{}{}{\color{RedViolet}\bfseries}{.}{ }{}
\theoremstyle{dotlessP}
\newtheorem{problem}{Problem}
\newtheoremstyle{dotlessS}{}{}{}{}{\color{ForestGreen}\bfseries}{.}{ }{}
\theoremstyle{dotlessS}
\newtheorem*{solution}{Solution}

\theoremstyle{remark}
\newtheorem*{remark}{Remark}
\theoremstyle{definition}
\newtheorem*{theorem}{Theorem}
\newtheorem*{lemma}{Lemma}

\tcbset{
  base/.style={
    arc=0mm,
    bottomtitle=0.5mm,
    boxrule=0mm,
    colbacktitle=black!10!white,
    coltitle=black,
    fonttitle=\bfseries,
    left=2.5mm,
    leftrule=1mm,
    right=3.5mm,
    title={#1},
    toptitle=0.75mm,
  }
}

\definecolor{brandblue}{rgb}{0.34, 0.7, 1}

% Key theorems (titled, blue).
\newtcolorbox{mainbox}[1]{
  enhanced jigsaw,
  breakable,
  colframe=brandblue,
  base={#1}
}

\newtcolorbox{subbox}[1]{
  colframe=black!30!white,
  base={#1}
}

% Key definitions and notation.
\newtcolorbox{definition_box}{
enhanced jigsaw,
breakable,
pad at break*=1mm,
colback=lightgray!12!white,
colframe=RedViolet!45!black,
colbacktitle=Gray!38!white,
before={\vspace{4mm}},
after={\vspace{2mm}},
sharp corners
}

% Exercises and open-ended problems.
\newtcolorbox{problem_box}{
enhanced jigsaw,
breakable,
pad at break*=1mm,
colback=white!5!white,
colframe=gray!75!black,
colbacktitle=white!28!white,
sharp corners,
before={\vspace{6mm}},
after={\vspace{2mm}}
}

\newtcolorbox{problem_box_1}{
enhanced jigsaw,
breakable,
pad at break*=1mm,
colback=SeaGreen!5!white,
colframe=SeaGreen!75!black,
colbacktitle=Black!28!white,
sharp corners,
before={\vspace{6mm}},
after={\vspace{2mm}}
}

\newtcolorbox{problem_box_title}[1]{
enhanced jigsaw,
breakable,
pad at break*=1mm,
colback=Apricot!5!white,
colframe=Melon!75!black,
sharp corners,
before={\vspace{2mm}},
after={\vspace{2mm}},
title={#1},
fonttitle=\bfseries
}

% Propositions, lemmas, proofs, and general text (unframed).
\newtcolorbox{claim_box}{
enhanced jigsaw,
breakable,
pad at break*=1mm,
colback=white,
colframe=white,
borderline west={2pt}{0pt}{white},
sharp corners,
before={\vspace{1mm}},
after={\vspace{0.8mm}}
}

% Examples.
\newtcolorbox{example_box}{
enhanced jigsaw,
breakable,
pad at break*=1mm,
colback=lightgray!12!white,
colframe=lightgray!15!white,
borderline west={2pt}{0pt}{black},
sharp corners,
before={\vspace{2mm}},
after={\vspace{2mm}}
}

% Cautions and remarks.
\newtcolorbox{remark_box}{
enhanced jigsaw,
breakable,
pad at break*=1mm,
colback=Apricot!5!white,
colframe=Melon!75!black,
sharp corners,
before={\vspace{2mm}},
after={\vspace{2mm}}
}


\begin{document}

%-----------------------------------------------------------------------------------
%   TOP MATTERS
%-----------------------------------------------------------------------------------


\title{Linear Transformations}
\author{Pratyay Sen}
\date{\today}

\maketitle

\tableofcontents

\clearpage


%-----------------------------------------------------------------------------------
%   MAIN BODY
%-----------------------------------------------------------------------------------

\section{Linear Maps}

\begin{definition_box}
\textbf{Definition:} Let $V$ and $W$ be vector spaces over the same
field $\mathbb{F}$. A map $T:V\longrightarrow W$ is a \emph{linear map},
or \emph{linear transformation}, if for all $u,v\in V$ and
$a\in\mathbb{F}$:
\begin{enumerate}
    \item (Additivity) $T(u+v)=T(u)+T(v)$;
    \item (Homogeneity) $T(av)=aT(v)$.
\end{enumerate}
Equivalently, for all $u,v\in V$ and $a,b\in\mathbb{F}$,
\[
T(au+bv)=aT(u)+bT(v).
\]
\end{definition_box}

\begin{claim_box}
\textbf{Proposition (linear maps send zero to zero):} If
$T:V\longrightarrow W$ is a linear map, then
\[
T(0_V)=0_W.
\]

\textit{Proof:} By additivity and $0_V+0_V=0_V$,
\[
T(0_V)=T(0_V+0_V)=T(0_V)+T(0_V).
\]
Adding the additive inverse of $T(0_V)$ to both sides in $W$ gives
$0_W=T(0_V)$, as required. \hfill$\square$
\end{claim_box}

\begin{claim_box}
\textbf{Proposition (linear maps on a one-dimensional space):} Let
$V$ be a one-dimensional vector space over $\mathbb{F}$ and let
$T:V\longrightarrow V$ be linear. There exists a unique scalar
$\lambda\in\mathbb{F}$ such that
\[
T(v)=\lambda v\qquad\text{for every }v\in V.
\]

\textit{Proof:} Choose a basis $\{e\}$ of $V$. Since $T(e)\in V$ and
$e$ spans $V$, there is a scalar $\lambda\in\mathbb{F}$ such that
\[
T(e)=\lambda e.
\]
Every $v\in V$ has the form $v=ae$ for some $a\in\mathbb{F}$.
By linearity and commutativity of scalar multiplication in the field,
\[
T(v)=T(ae)=aT(e)=a(\lambda e)=\lambda(ae)=\lambda v.
\]
Thus the same scalar $\lambda$ works for every vector in $V$.

For uniqueness, suppose also that $T(v)=\mu v$ for every $v\in V$.
Evaluating at $e$ gives $(\lambda-\mu)e=0_V$. Since $e\neq0_V$,
this implies $\lambda-\mu=0$, so $\lambda=\mu$.
\hfill$\square$
\end{claim_box}

\begin{claim_box}
\textbf{Note:} Conversely, multiplication by any fixed scalar is
a linear map. Hence these are exactly the linear maps from a
one-dimensional vector space to itself, including the zero map when
$\lambda=0$.
\end{claim_box}

\begin{definition_box}
\textbf{Notation:} The set of all linear maps from $V$ to $W$ over
$\mathbb{F}$ is denoted by $\mathcal{L}(V,W)$:
\[
\mathcal{L}(V,W)=\{T:V\longrightarrow W:T\text{ is linear}\}.
\]
Thus, to say that $T$ is a linear map from $V$ to $W$, we write
$T\in\mathcal{L}(V,W)$. The symbol $\mathcal{L}(V,W)$ denotes the whole set
of such maps, not a single map. When $W=V$, we abbreviate
$\mathcal{L}(V,V)$ to $\mathcal{L}(V)$; its elements are called
\emph{linear operators on $V$}.
\end{definition_box}

\begin{mainbox}{Vector space of linear maps}
Let $V$ and $W$
be vector spaces over the same field $\mathbb{F}$. Then
$\mathcal{L}(V,W)$ is a vector space over $\mathbb{F}$ under pointwise
addition and scalar multiplication, defined by
\[
(S+T)(v)=S(v)+T(v),
\qquad (aT)(v)=aT(v)
\]
for $S,T\in\mathcal{L}(V,W)$, $a\in\mathbb{F}$, and $v\in V$.
Here the operations on the right take place in $W$.
\end{mainbox}

\begin{claim_box}
\textit{Proof:} First we check closure. Let $S,T\in\mathcal{L}(V,W)$,
$u,v\in V$, and $a,b,c\in\mathbb{F}$. By linearity of $S$ and $T$,
\[
\begin{aligned}
(S+T)(au+bv)
&=S(au+bv)+T(au+bv)\\
&=aS(u)+bS(v)+aT(u)+bT(v)\\
&=a(S+T)(u)+b(S+T)(v).
\end{aligned}
\]
Thus $S+T$ is linear. Similarly,
\[
\begin{aligned}
(cT)(au+bv)
&=c\bigl(aT(u)+bT(v)\bigr)\\
&=a(cT)(u)+b(cT)(v),
\end{aligned}
\]
so $cT$ is linear as well.

The zero map $0_{\mathcal{L}}:V\longrightarrow W$, defined by
$0_{\mathcal{L}}(v)=0_W$ for every $v\in V$, is linear and satisfies
$T+0_{\mathcal{L}}=T$. For each $T$, the map $-T=(-1)T$ is linear and
\[
(T+(-T))(v)=T(v)-T(v)=0_W,
\]
so it is the additive inverse of $T$.

The remaining axioms hold pointwise because they hold in $W$. Explicitly,
for $R,S,T\in\mathcal{L}(V,W)$, $a,b\in\mathbb{F}$, and every $v\in V$,
\[
\begin{aligned}
(S+T)(v)&=S(v)+T(v)=T(v)+S(v)=(T+S)(v),\\
((R+S)+T)(v)&=(R(v)+S(v))+T(v)\\
&=R(v)+(S(v)+T(v))=(R+(S+T))(v),\\
(a(S+T))(v)&=a(S(v)+T(v))=(aS+aT)(v),\\
((a+b)T)(v)&=(a+b)T(v)=(aT+bT)(v),\\
(a(bT))(v)&=a(bT(v))=(ab)T(v)=((ab)T)(v),\\
(1T)(v)&=T(v).
\end{aligned}
\]
Two maps are equal when they agree on every vector in their domain.
Consequently, these pointwise identities establish all the corresponding
vector-space axioms for the maps themselves. Hence $\mathcal{L}(V,W)$ is
a vector space over $\mathbb{F}$. \hfill$\square$
\end{claim_box}

\begin{claim_box}
\textbf{Note:} The vectors of this space are linear maps.
Scalar multiplication here means multiplying the output of a map by a
scalar; it is not composition of maps. No finite-dimensional assumption
on $V$ or $W$ is needed.
\end{claim_box}

\begin{definition_box}
\textbf{Definition (product of linear maps):} Let $U,V,W$ be vector
spaces over the same field $\mathbb{F}$. If $T\in\mathcal{L}(U,V)$ and
$S\in\mathcal{L}(V,W)$, their \emph{product} $ST$ is the composition
\[
ST=S\circ T:U\longrightarrow W,
\qquad (ST)(u)=S(T(u)).
\]
The map on the right acts first: apply $T$, then $S$. The intermediate
spaces must match so that the output of $T$ can be used as input to $S$.
This product is not pointwise multiplication of the outputs.
\end{definition_box}

\begin{claim_box}
\textbf{Proposition:} The product of two composable linear maps is
linear. In the notation above, $ST\in\mathcal{L}(U,W)$.

\textit{Proof:} For $u,v\in U$ and $a,b\in\mathbb{F}$,
\[
\begin{aligned}
(ST)(au+bv)&=S(T(au+bv))\\
&=S(aT(u)+bT(v))\\
&=aS(T(u))+bS(T(v))\\
&=a(ST)(u)+b(ST)(v).
\end{aligned}
\]
Thus $ST$ is linear. \hfill$\square$
\end{claim_box}

\begin{claim_box}
\textbf{Properties:} Whenever the domains and codomains match,
composition is associative and distributes over addition:
\[
(RS)T=R(ST),\qquad
S(T_1+T_2)=ST_1+ST_2,\qquad
(S_1+S_2)T=S_1T+S_2T.
\]
Also, $a(ST)=(aS)T=S(aT)$ for $a\in\mathbb{F}$.
The identity map $I_V:V\longrightarrow V$, defined by $I_V(v)=v$, is
linear. For $T\in\mathcal{L}(U,V)$, it satisfies
\[
I_V T=T=T I_U.
\]
In particular, two operators in $\mathcal{L}(V)$ can always be multiplied,
but their product need not commute.
\end{claim_box}

\begin{example_box}
\textbf{Example:} Define operators on $\mathbb{R}^2$ by
\[
T(x,y)=(x,0),\qquad S(x,y)=(y,x).
\]
Then
\[
(ST)(x,y)=(0,x),\qquad (TS)(x,y)=(y,0).
\]
For instance, $(ST)(1,0)=(0,1)$ but $(TS)(1,0)=(0,0)$, so $ST\neq TS$.
For maps between different spaces, the reverse product may not even be
defined.
\end{example_box}

\begin{claim_box}
\textbf{Proposition (images of lines):} Let
$L:\mathbb{R}^n\longrightarrow\mathbb{R}^m$ be a linear map. The image of
every line in $\mathbb{R}^n$ is either a line or a single point in
$\mathbb{R}^m$.

\textit{Proof:} Any line $\ell$ in $\mathbb{R}^n$ can be written as
\[
\ell=\{p+tv:t\in\mathbb{R}\},
\]
where $p\in\mathbb{R}^n$ is a point on the line and
$v\in\mathbb{R}^n\setminus\{0\}$ is its direction vector. By linearity,
\[
L(p+tv)=L(p)+tL(v)
\qquad\text{for every }t\in\mathbb{R}.
\]
Therefore the image of the entire line is
\[
L(\ell)=\{L(p)+tL(v):t\in\mathbb{R}\}.
\]
There are two cases:
\begin{enumerate}
    \item If $L(v)\neq0$, this set is the line through $L(p)$ with
    direction vector $L(v)$.
    \item If $L(v)=0$, then $L(p+tv)=L(p)$ for every $t$, so
    $L(\ell)=\{L(p)\}$ is a single point.
\end{enumerate}
Thus a linear map sends every line to a line or a point. The original line
need not pass through the origin, and the image point need not be the
origin. \hfill$\square$
\end{claim_box}

\begin{problem_box}
\textbf{Open-ended problem (a converse for maps of the plane):} Let
$f:\mathbb{R}^2\longrightarrow\mathbb{R}^2$ be a map satisfying:
\begin{enumerate}
    \item $f$ is injective;
    \item $f(0)=0$;
    \item for every line $\ell\subseteq\mathbb{R}^2$, its image
    $f(\ell)=\{f(x):x\in\ell\}$ is a line in $\mathbb{R}^2$.
\end{enumerate}
Prove that $f$ is a linear map over $\mathbb{R}$; that is, prove that
\[
f(au+bv)=af(u)+bf(v)
\]
for all $u,v\in\mathbb{R}^2$ and $a,b\in\mathbb{R}$.
Here the third assumption says that the image is an entire line, not merely
a subset of a line. No continuity assumption is given.
\end{problem_box}

\begin{claim_box}
\textbf{Differentiation:} Let $\mathcal{P}_n(\mathbb{F})$
denote the polynomials of degree at most $n$ over $\mathbb{F}$, including
the zero polynomial. For $n\geq1$, formal differentiation defines a map
\[
D:\mathcal{P}_n(\mathbb{F})\longrightarrow\mathcal{P}_{n-1}(\mathbb{F}),
\qquad D(p)=p'.
\]
Explicitly,
\[
D(a_0+a_1x+\cdots+a_n x^n)
=a_1+2a_2x+\cdots+na_n x^{n-1}.
\]
For any $p,q\in\mathcal{P}_n(\mathbb{F})$ and $a,b\in\mathbb{F}$,
\[
D(ap+bq)=aD(p)+bD(q),
\]
so $D$ is linear. Over $\mathbb{R}$, for example,
$D(3x^2+2x+5)=6x+2$. A nonzero constant polynomial is sent to zero,
showing that a linear map need not be injective.
\end{claim_box}

\begin{claim_box}
\textbf{Integration:} Let $C([0,1],\mathbb{R})$ denote
the vector space of continuous real-valued functions on $[0,1]$, with
pointwise addition and scalar multiplication. Define
\[
J:C([0,1],\mathbb{R})\longrightarrow C([0,1],\mathbb{R}),
\qquad (Jf)(x)=\int_0^x f(t)\,dt.
\]
The function $Jf$ is continuous, so the map has the stated codomain.
For any $f,g\in C([0,1],\mathbb{R})$ and $a,b\in\mathbb{R}$,
\[
\begin{aligned}
(J(af+bg))(x)
&=\int_0^x \bigl(af(t)+bg(t)\bigr)\,dt\\
&=a\int_0^x f(t)\,dt+b\int_0^x g(t)\,dt\\
&=a(Jf)(x)+b(Jg)(x).
\end{aligned}
\]
Thus $J$ is linear. For example, if $f(t)=t^2$, then $(Jf)(x)=x^3/3$.
The fixed lower limit specifies a unique antiderivative with value zero
at $0$; an indefinite integral with an unspecified constant does not
specify a single function-valued map.

Integration over a fixed interval also gives a linear map with scalar
values:
\[
I:C([0,1],\mathbb{R})\longrightarrow\mathbb{R},
\qquad I(f)=\int_0^1 f(t)\,dt.
\]
Indeed, $I(af+bg)=aI(f)+bI(g)$.
\end{claim_box}

\begin{claim_box}
\textbf{Matrix maps:} Fix an $m\times n$ matrix
$A=(a_{ij})$ with entries in $\mathbb{F}$. Regarding vectors as column
vectors, define
\[
T_A:\mathbb{F}^n\longrightarrow\mathbb{F}^m,
\qquad T_A(x)=Ax.
\]
The $i$th coordinate of $T_A(x)$ is $\sum_{j=1}^n a_{ij}x_j$.
For $u,v\in\mathbb{F}^n$ and $a,b\in\mathbb{F}$,
\[
T_A(au+bv)=A(au+bv)=aAu+bAv=aT_A(u)+bT_A(v),
\]
so $T_A$ is linear. For example,
\[
A=\begin{pmatrix}1&2&0\\0&-1&1\end{pmatrix}
\quad\text{gives}\quad
T_A:\mathbb{F}^3\longrightarrow\mathbb{F}^2,
\qquad T_A(x,y,z)=(x+2y,-y+z).
\]
\end{claim_box}

\begin{example_box}
\textbf{Example:} The map
\[
T:\mathbb{R}^2\longrightarrow\mathbb{R}^2,
\qquad T(x,y)=(2x,3y),
\]
is linear: it preserves both vector addition and scalar multiplication.
\end{example_box}

\begin{claim_box}
\textbf{Counterexample (homogeneity without additivity):} Regard
$\mathbb{R}^2$ and $\mathbb{R}$ as vector spaces over $\mathbb{R}$. Define
\[
H:\mathbb{R}^2\longrightarrow\mathbb{R},
\qquad
H(x,y)=
\begin{cases}
\dfrac{x^3}{x^2+y^2}, & (x,y)\neq(0,0),\\[4pt]
0, & (x,y)=(0,0).
\end{cases}
\]
For every nonzero $a\in\mathbb{R}$ and $(x,y)\neq(0,0)$,
\[
H(ax,ay)=\frac{a^3x^3}{a^2(x^2+y^2)}=aH(x,y).
\]
The same identity holds when $a=0$ or $(x,y)=(0,0)$, since both sides
are zero. Thus $H$ is homogeneous for every real scalar, including
negative scalars. However,
\[
H((1,0)+(0,1))=H(1,1)=\frac12,
\qquad H(1,0)+H(0,1)=1+0=1.
\]
So $H$ is not additive and therefore is not linear.
\end{claim_box}

\begin{claim_box}
\textbf{Counterexample (additivity without homogeneity):} Regard
$\mathbb{C}$ as a vector space over $\mathbb{C}$ and define complex
conjugation by
\[
K:\mathbb{C}\longrightarrow\mathbb{C},
\qquad K(z)=\overline{z}.
\]
It is additive because, for all $z,w\in\mathbb{C}$,
\[
K(z+w)=\overline{z+w}=\overline{z}+\overline{w}=K(z)+K(w).
\]
But homogeneity fails for the complex scalar $i$ and the vector $1$:
\[
K(i\cdot1)=\overline{i}=-i,
\qquad iK(1)=i.
\]
Hence $K$ is not linear over $\mathbb{C}$.
\end{claim_box}

\begin{remark_box}
\textbf{Note (the scalar field matters):} Conjugation does satisfy
$K(az)=aK(z)$ for every real scalar $a$. Thus it is linear when
$\mathbb{C}$ is regarded as a vector space over $\mathbb{R}$, but not when
the scalar field is $\mathbb{C}$. These counterexamples show why both
additivity and homogeneity are required in the definition of a linear map.
\end{remark_box}

\begin{claim_box}
\textbf{Definition (graph of a linear map):} Let
$T:V\longrightarrow W$ be linear over $\mathbb{F}$. The \emph{graph}
of $T$ is
\[
\Gamma(T)=\{(v,T(v)):v\in V\}\subseteq V\times W.
\]
Here $V\times W$ is a vector space with componentwise operations:
\[
(v,w)+(v',w')=(v+v',w+w'),\qquad a(v,w)=(av,aw).
\]
Unlike the range, the graph records both the input and its output.
\end{claim_box}

\begin{claim_box}
\textbf{Proposition (graphs and linearity):} A map
$f:V\longrightarrow W$ is linear if and only if its graph
$\Gamma(f)$ is a vector subspace of $V\times W$.

\textit{Proof:} If $f$ is linear, then $(0_V,0_W)$ belongs to its
graph. For $u,v\in V$ and $a,b\in\mathbb{F}$,
\[
a(u,f(u))+b(v,f(v))=(au+bv,af(u)+bf(v))
=(au+bv,f(au+bv)),
\]
which also belongs to the graph. Thus the graph is a subspace.

Conversely, suppose $\Gamma(f)$ is a subspace. The same linear
combination $(au+bv,af(u)+bf(v))$ belongs to $\Gamma(f)$.
The unique pair in the graph with first component $au+bv$ has second
component $f(au+bv)$. Hence
$f(au+bv)=af(u)+bf(v)$, proving linearity. \hfill$\square$
\end{claim_box}

\begin{claim_box}
\textbf{Proposition (the graph is isomorphic to the domain):}
For a linear map $T:V\longrightarrow W$, the map
\[
G:V\longrightarrow\Gamma(T),\qquad G(v)=(v,T(v)),
\]
is an isomorphism. In particular, if $V$ is finite-dimensional, then
$\dim\Gamma(T)=\dim V$.

\textit{Proof:} Linearity follows from the componentwise operations
and the linearity of $T$. Equality $G(u)=G(v)$ forces $u=v$ by
comparing first components, so $G$ is injective. Every element of
$\Gamma(T)$ has the form $G(v)$, so $G$ is surjective. Its inverse
is the first-coordinate projection $(v,T(v))\mapsto v$.
The dimension statement follows because isomorphisms preserve
dimension. \hfill$\square$
\end{claim_box}

\begin{claim_box}
\textbf{Proposition (recovering the kernel and range):} The
second-coordinate projection restricted to the graph is the linear map
\[
Q:\Gamma(T)\longrightarrow W,\qquad Q(v,T(v))=T(v).
\]
It satisfies
\[
\ker Q=\{(v,0_W):v\in\ker T\},\qquad
\operatorname{range}Q=\operatorname{range}T.
\]
Equivalently,
\[
\Gamma(T)\cap(V\times\{0_W\})=(\ker T)\times\{0_W\}.
\]
Thus $T$ is injective exactly when this intersection is
$\{(0_V,0_W)\}$, and $T$ is surjective exactly when $Q$ is surjective.

\textit{Proof:} Projection onto the second component is linear.
A pair $(v,T(v))$ is in its kernel precisely when $T(v)=0_W$,
and its possible outputs are precisely the values attained by $T$.
These observations prove both identities and the intersection formula.
The final assertions follow from the kernel criterion for injectivity
and the definition of surjectivity. \hfill$\square$
\end{claim_box}

\begin{claim_box}
\textbf{Proposition (the graph as a direct complement):} We have
\[
V\times W=\Gamma(T)\oplus(\{0_V\}\times W).
\]
Here the direct sum means that each pair is uniquely the sum of one
vector in each of these two subspaces.

\textit{Proof:} Every $(v,w)\in V\times W$ can be written as
\[
(v,w)=(v,T(v))+(0_V,w-T(v)).
\]
The first summand belongs to the graph and the second to
$\{0_V\}\times W$. If a graph element $(v,T(v))$ also lies in
$\{0_V\}\times W$, then $v=0_V$ and $T(v)=0_W$. The intersection
is therefore the zero subspace, making the decomposition unique.
\hfill$\square$
\end{claim_box}

\begin{example_box}
\textbf{Example:} For $T:\mathbb{R}^2\longrightarrow\mathbb{R}$,
$T(x,y)=x+2y$, identify $\mathbb{R}^2\times\mathbb{R}$ with
$\mathbb{R}^3$. Then
\[
\Gamma(T)=\{(x,y,x+2y):x,y\in\mathbb{R}\}
=\operatorname{span}\{(1,0,1),(0,1,2)\}.
\]
The graph is the plane $z=x+2y$, of dimension $2$, while the range is
$\mathbb{R}$, of dimension $1$. Its intersection with $z=0$ is the
line $\{(-2t,t,0):t\in\mathbb{R}\}$, corresponding to
$\ker T=\operatorname{span}\{(-2,1)\}$.
\end{example_box}

\subsection{Isomorphisms}

\begin{definition_box}
\textbf{Definition (isomorphism):} Let $V$ and $W$ be vector spaces
over the same field $\mathbb{F}$. A map $T:V\longrightarrow W$ is a
\emph{vector-space isomorphism} if it is linear and bijective. Thus every
$w\in W$ has exactly one preimage $v\in V$ under $T$.
If an isomorphism exists, we call $V$ and $W$ \emph{isomorphic} and write
\[
V\cong W.
\]
This means that the spaces have the same vector-space structure, even if
their elements are different kinds of objects.
\end{definition_box}

\begin{definition_box}
\textbf{Definition (invertibility):} A linear map
$T:V\longrightarrow W$ is \emph{invertible} if there is a map
$R:W\longrightarrow V$ such that
\[
RT=I_V,\qquad TR=I_W.
\]
Here $I_V$ and $I_W$ are the identity maps. These equations say that
$R(T(v))=v$ for every $v\in V$ and $T(R(w))=w$ for every $w\in W$:
$R$ undoes $T$ in both directions. The map $R$ is called the
\emph{inverse} of $T$ and is denoted by $T^{-1}$.
\end{definition_box}

\begin{claim_box}
The inverse is unique. Indeed, if $R$ and $Q$ are both inverses of $T$,
then associativity of composition gives
\[
R=R I_W=R(TQ)=(RT)Q=I_V Q=Q.
\]
\end{claim_box}

\begin{claim_box}
\textbf{Lemma (the inverse of an isomorphism is linear).} If
$T:V\longrightarrow W$ is an isomorphism, then its inverse
$T^{-1}:W\longrightarrow V$ is also an isomorphism.

\textit{Proof:} Since $T$ is bijective, its inverse exists and is bijective.
It remains to prove that $T^{-1}$ is linear. Let $w_1,w_2\in W$ and
$a,b\in\mathbb{F}$. Set
\[
v_1=T^{-1}(w_1),\qquad v_2=T^{-1}(w_2).
\]
Then $T(v_1)=w_1$ and $T(v_2)=w_2$. By linearity of $T$,
\[
T(av_1+bv_2)=aT(v_1)+bT(v_2)=aw_1+bw_2.
\]
Applying $T^{-1}$ to both sides gives
\[
\begin{aligned}
T^{-1}(aw_1+bw_2)
&=av_1+bv_2\\
&=aT^{-1}(w_1)+bT^{-1}(w_2).
\end{aligned}
\]
Thus $T^{-1}$ is linear and bijective, hence an isomorphism.
\hfill$\square$
\end{claim_box}

\begin{claim_box}
\textbf{Proposition (invertibility and isomorphisms):} A linear
map $T:V\longrightarrow W$ is invertible if and only if it is an
isomorphism. In that case $T^{-1}$ is automatically linear.

\textit{Proof:} Suppose $R$ is an inverse of $T$. If $T(u)=T(v)$,
applying $R$ gives $u=v$, so $T$ is injective. For every $w\in W$,
the identity $T(R(w))=w$ shows that $T$ is surjective. Thus $T$ is
linear and bijective, hence an isomorphism.
Conversely, an isomorphism is bijective, so each $w\in W$ has a unique
preimage. Assigning that preimage to $w$ defines a map $T^{-1}$
satisfying both inverse identities. The preceding lemma proves that
this inverse is linear. \hfill$\square$
\end{claim_box}

\begin{claim_box}
\textbf{Interpretation (solving equations):} Invertibility means
that for every $w\in W$, the equation $T(v)=w$ has exactly one solution,
namely $v=T^{-1}(w)$. Existence for every right-hand side is surjectivity;
uniqueness is injectivity. Both are required.
\end{claim_box}

\begin{example_box}
\textbf{Example (scalar multiplication):} On a nonzero vector
space $V$, the map $T(v)=\lambda v$ is invertible exactly when
$\lambda\neq0$. In that case $T^{-1}(v)=\lambda^{-1}v$, since composing
the two maps in either order gives $v$. If $\lambda=0$, all vectors
have the same image, so the map is not injective. In particular,
$T(x,y)=(2x,3y)$ on $\mathbb{R}^2$ is invertible with
$T^{-1}(x,y)=(x/2,y/3)$.
\end{example_box}

\begin{example_box}
\textbf{Example (a coefficient isomorphism):} Let
$\mathcal{P}_2(\mathbb{F})$ be the space of polynomials of degree at most
$2$, including zero. The map
\[
T:\mathcal{P}_2(\mathbb{F})\longrightarrow\mathbb{F}^3,
\qquad T(a+bx+cx^2)=(a,b,c),
\]
is linear and bijective. Its inverse sends $(a,b,c)$ to $a+bx+cx^2$.
Consequently, $\mathcal{P}_2(\mathbb{F})\cong\mathbb{F}^3$.
\end{example_box}

\begin{claim_box}
\textbf{Intuition (a reversible change of description):} An
isomorphism lets us describe every vector in another space without losing
information or changing how linear combinations work:
\[
T(a_1v_1+\cdots+a_k v_k)=a_1T(v_1)+\cdots+a_k T(v_k).
\]
Injectivity means that different vectors remain distinguishable, and
surjectivity means that every vector in the target is represented.
The inverse lets us recover the original vector. This does not mean the
spaces are literally equal, nor does it require preservation of lengths
or angles.
\end{claim_box}

\begin{example_box}
\textbf{Example (stretching and shearing):} The map
\[
S:\mathbb{R}^2\longrightarrow\mathbb{R}^2,
\qquad S(x,y)=(2x+y,y)
\]
is linear. Given $(u,v)\in\mathbb{R}^2$, its unique preimage is
$((u-v)/2,v)$, so
\[
S^{-1}(u,v)=\left(\frac{u-v}{2},v\right).
\]
Thus $S$ is an isomorphism. It sends $(1,0)$ to $(2,0)$, showing that an
isomorphism need not preserve Euclidean length. An isomorphism from a
vector space to itself is called an \emph{automorphism}.
\end{example_box}

\begin{example_box}
\textbf{Example (complex numbers as real vectors):} Regard
$\mathbb{C}$ as a vector space over $\mathbb{R}$. Then
\[
\Phi:\mathbb{C}\longrightarrow\mathbb{R}^2,
\qquad \Phi(a+bi)=(a,b)
\]
is a real-linear isomorphism, with inverse $(a,b)\mapsto a+bi$.
The scalar field is essential: $\mathbb{C}$ has dimension $2$ over
$\mathbb{R}$ but dimension $1$ over $\mathbb{C}$.
\end{example_box}

\begin{example_box}
\textbf{Nonexamples (why all conditions matter):}
\begin{enumerate}
    \item The projection $P:\mathbb{R}^2\longrightarrow\mathbb{R}$,
    $P(x,y)=x$, is linear and surjective, but not injective:
    $P(0,0)=P(0,1)$. Hence it is not an isomorphism.
    \item The inclusion $J:\mathbb{R}\longrightarrow\mathbb{R}^2$,
    $J(t)=(t,0)$, is linear and injective, but not surjective: $(0,1)$
    is not in its image. With the horizontal axis as its codomain instead,
    $J$ would be an isomorphism.
    \item The map $f:\mathbb{R}\longrightarrow\mathbb{R}$,
    $f(x)=x^3$, is bijective but not linear, since
    $f(1+1)=8\neq2=f(1)+f(1)$. A bijection alone is not enough.
\end{enumerate}
\end{example_box}

\begin{claim_box}
\textbf{Proposition (isomorphisms carry bases to bases):} If
$T:V\longrightarrow W$ is an isomorphism and $B$ is a basis of $V$,
then $T(B)=\{T(v):v\in B\}$ is a basis of $W$.

\textit{Proof:} Every $w\in W$ equals $T(v)$ for some $v\in V$.
Expressing $v$ as a finite linear combination of elements of $B$ and
applying $T$ shows that $T(B)$ spans $W$.
For independence, suppose $v_1,\ldots,v_k$ are distinct elements of $B$ and
\[
a_1T(v_1)+\cdots+a_k T(v_k)=0_W.
\]
By linearity, $T(a_1v_1+\cdots+a_k v_k)=T(0_V)$. Injectivity implies
$a_1v_1+\cdots+a_k v_k=0_V$. Independence of $B$ then gives
$a_1=\cdots=a_k=0$. Thus $T(B)$ is a basis of $W$.
\hfill$\square$
\end{claim_box}

\begin{example_box}
\textbf{Example (coordinates):} If $v_1,\ldots,v_n$ is an ordered
basis of $V$, the coordinate map
\[
\sum_{i=1}^n a_i v_i\longmapsto(a_1,\ldots,a_n)
\]
is an isomorphism $V\cong\mathbb{F}^n$. Different choices of basis
generally give different coordinate maps. In particular,
$\mathcal{P}_n(\mathbb{F})\cong\mathbb{F}^{n+1}$ using the basis
$1,x,\ldots,x^n$.
\end{example_box}

\begin{remark_box}
\textbf{Caution (one-sided inverses):} A map $R:W\longrightarrow V$
with $RT=I_V$ is a \emph{left inverse} of $T$, and this identity alone
implies injectivity. A map with $TR=I_W$ is a \emph{right inverse},
and this identity alone implies surjectivity. In general, one identity
does not imply the other. For example, let
\[
J:\mathbb{R}\longrightarrow\mathbb{R}^2,\quad J(t)=(t,0),
\qquad P:\mathbb{R}^2\longrightarrow\mathbb{R},\quad P(x,y)=x.
\]
Then $PJ=I_{\mathbb{R}}$, but $JP(x,y)=(x,0)$ is not
$I_{\mathbb{R}^2}(x,y)$. Neither $J$ nor $P$ is invertible.
When the domain and codomain have the same finite dimension, however,
the \hyperref[thm:dimension-criteria]{dimension theorem above} shows that
either inverse identity guarantees
invertibility, and the one-sided inverse is then the unique inverse.
\end{remark_box}

\begin{claim_box}
\textbf{Proposition (composition and isomorphism):} If
$T:U\longrightarrow V$ and $S:V\longrightarrow W$ are isomorphisms, then
$ST:U\longrightarrow W$ is an isomorphism, and
\[
{(ST)}^{-1}=T^{-1}S^{-1}.
\]

\textit{Proof:} Composition preserves linearity. Moreover,
\[
(T^{-1}S^{-1})(ST)=I_U,
\qquad (ST)(T^{-1}S^{-1})=I_W,
\]
so $ST$ is bijective with the stated inverse. \hfill$\square$

Together with the identity map and the inverse lemma, this shows that
being isomorphic is reflexive, symmetric, and transitive: $V\cong V$;
if $V\cong W$, then $W\cong V$; and if $U\cong V$ and $V\cong W$,
then $U\cong W$.
\end{claim_box}

\begin{mainbox}{Theorem (dimension of the space of linear maps)}
Let
$V$ and $W$ be finite-dimensional vector spaces over the same field
$\mathbb{F}$. Then
\[
\dim\mathcal{L}(V,W)=(\dim V)(\dim W).
\]
\end{mainbox}

\begin{claim_box}
\textit{Proof:} Put $n=\dim V$ and $m=\dim W$. If either space is
the zero space, the only linear map from $V$ to $W$ is the zero map.
Thus $\dim\mathcal{L}(V,W)=0=mn$ in this case.

Otherwise, fix ordered bases $\mathcal{B}=(v_1,\ldots,v_n)$ of $V$
and $\mathcal{C}=(w_1,\ldots,w_m)$ of $W$. We use the matrix
representation and its linearity, developed in the later subsection
\emph{Matrices}. The matrix representation map
\[
M:\mathcal{L}(V,W)\longrightarrow M_{m\times n}(\mathbb{F})
\]
is linear. It is also bijective: a matrix $A=(a_{ij})$ specifies
exactly one linear map, namely
\[
T_A\left(\sum_{j=1}^n x_j v_j\right)
=\sum_{i=1}^m\left(\sum_{j=1}^n a_{ij}x_j\right)w_i.
\]
Unique coordinates in $\mathcal{B}$ make this well-defined, and the
formula is linear in the $x_j$. Its representing matrix is $A$.
Conversely, any linear map with matrix $A$ must have this formula,
because its values on the basis vectors determine its values on every
vector. Hence $M$ is an isomorphism.

For $1\leq i\leq m$ and $1\leq j\leq n$, let $E_{ij}$ be the
$m\times n$ matrix with a $1$ in row $i$, column $j$, and zeros
elsewhere. Every matrix has the unique expansion
\[
A=\sum_{i=1}^m\sum_{j=1}^n a_{ij}E_{ij}.
\]
Indeed, the entry in row $i$, column $j$ of this sum is $a_{ij}$.
Thus the $mn$ matrices $E_{ij}$ span the matrix space and are linearly
independent, so they form a basis. Since an isomorphism carries a basis
to a basis, their inverse images under $M$ form a basis of
$\mathcal{L}(V,W)$. Consequently,
\[
\dim\mathcal{L}(V,W)=mn=(\dim V)(\dim W).
\]
\hfill$\square$
\end{claim_box}

\subsection{Null spaces and ranges}

\begin{definition_box}
\textbf{Definition (null space or kernel):} Let
$T:V\longrightarrow W$ be a linear map over $\mathbb{F}$. The
\emph{null space} or \emph{kernel} of $T$ is
\[
\operatorname{null}T=\ker T=\{v\in V:T(v)=0_W\}.
\]
It consists of all input vectors that $T$ sends to zero. In particular,
the null space is a subset of the \emph{domain} $V$, not the codomain $W$.
It always contains $0_V$, but may also contain nonzero vectors.
\end{definition_box}

\begin{claim_box}
\textbf{Proposition:} The null space of $T$ is a subspace of $V$.

\textit{Proof:} Since $T(0_V)=0_W$, we have $0_V\in\ker T$.
If $u,v\in\ker T$ and $a,b\in\mathbb{F}$, linearity gives
\[
T(au+bv)=aT(u)+bT(v)=a0_W+b0_W=0_W.
\]
Thus $au+bv\in\ker T$, proving that $\ker T$ is a subspace.
\hfill$\square$
\end{claim_box}

\begin{definition_box}
\textbf{Definition (range or image):} The \emph{range} or
\emph{image} of a linear map $T:V\longrightarrow W$ is
\[
\operatorname{range}T=\operatorname{im}T=T(V)
=\{T(v):v\in V\}\subseteq W.
\]
It consists of all output vectors that $T$ actually attains. It need not
equal the whole codomain $W$; equality holds exactly when $T$ is surjective.
\end{definition_box}

\begin{claim_box}
\textbf{Proposition:} The range of $T$ is a subspace of $W$.

\textit{Proof:} Since $T(0_V)=0_W$, the range contains $0_W$.
Let $w_1,w_2\in\operatorname{range}T$ and $a,b\in\mathbb{F}$.
By the definition of the range, there exist $v_1,v_2\in V$ such that
$w_1=T(v_1)$ and $w_2=T(v_2)$. By linearity,
\[
aw_1+bw_2=aT(v_1)+bT(v_2)=T(av_1+bv_2).
\]
Because $av_1+bv_2\in V$, this output belongs to
$\operatorname{range}T$. Hence the range contains zero and is closed under
linear combinations, so it is a subspace of $W$. \hfill$\square$
\end{claim_box}

\begin{claim_box}
\textbf{Summary:} The two subspaces live in different spaces:
\[
\ker T\text{ is a subspace of }V,
\qquad
\operatorname{range}T\text{ is a subspace of }W.
\]
\end{claim_box}

\begin{definition_box}
\textbf{Definition (rank):} The \emph{rank} of a linear map
$T:V\longrightarrow W$ is the dimension of its range:
\[
\operatorname{rank}T=\dim(\operatorname{range}T).
\]
Thus rank measures the number of linearly independent output directions
attained by $T$. If $V$ and $W$ are finite-dimensional, then
\[
0\leq\operatorname{rank}T\leq\min\{\dim V,\dim W\}.
\]

For a matrix $A\in\mathbb{F}^{m\times n}$, its \emph{rank} is the rank
of the map $T_A(x)=Ax$. Since the range of $T_A$ is the span of the
columns $A_1,\ldots,A_n$,
\[
\operatorname{rank}A
=\dim\operatorname{span}\{A_1,\ldots,A_n\}.
\]
Equivalently, it is the largest number of linearly independent columns
of $A$.
\end{definition_box}

\begin{example_box}
\textbf{Example (rank):} The projection
$P:\mathbb{R}^2\longrightarrow\mathbb{R}$, $P(x,y)=x$, has range
$\mathbb{R}$, so $\operatorname{rank}P=1$. The zero map has rank $0$,
and the identity map on $\mathbb{F}^n$ has rank $n$.
\end{example_box}

\begin{example_box}
\textbf{Example (projection):} For
\[
P:\mathbb{R}^2\longrightarrow\mathbb{R},\qquad P(x,y)=x,
\]
the equation $P(x,y)=0$ is equivalent to $x=0$. Hence
\[
\ker P=\{(0,y):y\in\mathbb{R}\}
=\operatorname{span}\{(0,1)\}.
\]
The null space is the vertical axis: the projection discards the entire
vertical component.
\end{example_box}

\begin{example_box}
\textbf{Example (differentiation):} Over $\mathbb{R}$, consider
\[
D:\mathcal{P}_n(\mathbb{R})\longrightarrow\mathcal{P}_{n-1}(\mathbb{R}),
\qquad D(p)=p',\quad n\geq1.
\]
A polynomial has zero derivative exactly when it is constant. Therefore
\[
\ker D=\{p:p(x)=c,\ c\in\mathbb{R}\}
=\operatorname{span}\{1\}.
\]
Here $1$ denotes the constant polynomial, and the zero polynomial is
included by taking $c=0$.
\end{example_box}

\begin{example_box}
\textbf{Example (matrix null space):} For the matrix map
\[
T_A:\mathbb{R}^3\longrightarrow\mathbb{R}^2,
\qquad T_A(x,y,z)=(x+2y,-y+z),
\]
finding the null space means solving the homogeneous system $Av=0$:
\[
x+2y=0,\qquad -y+z=0.
\]
Taking $y=t$ gives $x=-2t$ and $z=t$, so
\[
\ker T_A=\{(-2t,t,t):t\in\mathbb{R}\}
=\operatorname{span}\{(-2,1,1)\}.
\]
\end{example_box}

\begin{example_box}
\textbf{Examples (the two extremes):} The identity map $I_V$ has
$\ker I_V=\{0_V\}$, whereas the zero map $Z:V\longrightarrow W$,
defined by $Z(v)=0_W$, has $\ker Z=V$.
\end{example_box}

\begin{claim_box}
\textbf{Proposition (null space and injectivity):} A linear map
$T:V\longrightarrow W$ is injective if and only if $\ker T=\{0_V\}$.

\textit{Proof:} If $T$ is injective and $v\in\ker T$, then
$T(v)=0_W=T(0_V)$ implies $v=0_V$. Conversely, suppose
$\ker T=\{0_V\}$. If $T(u)=T(v)$, then
\[
T(u-v)=T(u)-T(v)=0_W,
\]
so $u-v\in\ker T$, giving $u=v$. Thus $T$ is injective.
\hfill$\square$
\end{claim_box}

\begin{mainbox}{Rank-Nullity Theorem}
Let $V$ be a
finite-dimensional vector space over $\mathbb{F}$, let $W$ be a vector
space over the same field, and let $T\in\mathcal{L}(V,W)$. Then
$\operatorname{range}T$ is finite-dimensional and
\[
\dim V=\dim(\ker T)+\dim(\operatorname{range}T).
\]
Note that $W$ need not be
finite-dimensional.
\end{mainbox}

\begin{claim_box}
\textit{Proof:} Since $\ker T$ is a subspace of the finite-dimensional
space $V$, it has a basis $u_1,\ldots,u_k$. Extend this to a basis
\[
u_1,\ldots,u_k,v_1,\ldots,v_r
\]
of $V$, so $\dim V=k+r$. We will show that
\[
T(v_1),\ldots,T(v_r)
\]
is a basis of $\operatorname{range}T$.

\emph{Spanning:} Let $w\in\operatorname{range}T$. There is some $v\in V$
with $T(v)=w$. Expanding $v$ in the chosen basis gives
\[
v=\sum_{i=1}^k a_i u_i+\sum_{j=1}^r b_j v_j.
\]
Because each $u_i$ belongs to $\ker T$, applying $T$ yields
\[
w=T(v)=\sum_{i=1}^k a_i T(u_i)+\sum_{j=1}^r b_j T(v_j)
=\sum_{j=1}^r b_j T(v_j).
\]
Thus $T(v_1),\ldots,T(v_r)$ spans the range.

\emph{Linear independence:} Suppose
\[
\sum_{j=1}^r c_j T(v_j)=0_W.
\]
By linearity, $T(\sum_{j=1}^r c_j v_j)=0_W$, so
$\sum_{j=1}^r c_j v_j\in\ker T$. Since the $u_i$ form a basis of the
kernel, there are scalars $d_1,\ldots,d_k$ such that
\[
\sum_{j=1}^r c_j v_j=\sum_{i=1}^k d_i u_i.
\]
Subtracting the right-hand side gives a linear relation among the basis
vectors of $V$. Hence all its coefficients vanish, and in particular
$c_1=\cdots=c_r=0$. This proves independence.

Consequently, the range has a finite basis of length $r$. Therefore
\[
\dim(\ker T)+\dim(\operatorname{range}T)=k+r=\dim V.
\]
Empty basis lists and empty sums are allowed, so the proof also covers
the zero map and the case $V=\{0_V\}$. \hfill$\square$
\end{claim_box}

\phantomsection%
\label{thm:dimension-criteria}%
\begin{mainbox}{Theorem (dimension criteria for linear maps)}
Let $V$ and
$W$ be finite-dimensional vector spaces over the same field $\mathbb{F}$.
Write $n=\dim V$ and $m=\dim W$. Then:
\begin{enumerate}
    \item If $n>m$, no linear map $T:V\longrightarrow W$ is injective.
    \item If $n<m$, no linear map $T:V\longrightarrow W$ is surjective.
    \item If $n=m$, a linear map $T:V\longrightarrow W$ is injective
    if and only if it is surjective, and either condition is equivalent
    to invertibility.
    \item $V$ and $W$ are isomorphic if and only if $n=m$.
\end{enumerate}
\end{mainbox}

\begin{claim_box}

\textit{Proof:} For any $T\in\mathcal{L}(V,W)$, put
$r=\dim(\operatorname{range}T)$ and $k=\dim(\ker T)$.
Rank-nullity and the fact that the range is a subspace of $W$ give
\[
n=k+r,\qquad k\geq0,\qquad r\leq m.
\]
Consequently, $r\leq\min\{n,m\}$. Moreover, the kernel criterion for
injectivity and the dimension criterion for equality of subspaces give
\[
\begin{aligned}
T\text{ is injective}&\ \Longleftrightarrow\ k=0
\ \Longleftrightarrow\ r=n,\\
T\text{ is surjective}&\ \Longleftrightarrow\
\operatorname{range}T=W\ \Longleftrightarrow\ r=m.
\end{aligned}
\]
These two equivalences prove the first three claims at once: if $n>m$,
then $r<n$, so injectivity is impossible; if $n<m$, then $r<m$, so
surjectivity is impossible; and if $n=m$, the two conditions coincide.
In the last case, either condition makes $T$ bijective and hence
invertible. They also show that an isomorphism must have $r=n=m$.

It remains to show that equal dimensions guarantee the existence of
an isomorphism. If $n=m$, choose bases $v_1,\ldots,v_n$ of $V$ and
$w_1,\ldots,w_n$ of $W$, and define
\[
T\left(\sum_{i=1}^n a_i v_i\right)=\sum_{i=1}^n a_i w_i.
\]
Unique basis expansions make $T$ well-defined, and coordinate addition
and scalar multiplication show that it is linear. Every vector of $W$
has the displayed form on the right, so $T$ is surjective. By the
equal-dimension case already proved, $T$ is invertible, hence an
isomorphism. Empty bases cover $n=m=0$. \hfill$\square$
\end{claim_box}

\begin{claim_box}
\textbf{Definition (the induced map on the quotient):} Let
$T:V\longrightarrow W$ be linear and put $U=\operatorname{null}T=\ker T$.
Recall that $V/U$ consists of cosets $v+U$, with operations
\[
(v+U)+(w+U)=(v+w)+U,\qquad a(v+U)=av+U.
\]
The \emph{map induced by $T$ on $V/\operatorname{null}T$} is
\[
\widetilde{T}:V/\operatorname{null}T\longrightarrow W,
\qquad \widetilde{T}(v+U)=T(v).
\]
Its input is a coset, and its output is the image under $T$ of any
representative of that coset. No finite-dimensional assumption is needed.
\end{claim_box}

\begin{claim_box}
\textbf{Proposition:} The induced map $\widetilde{T}$ is
well-defined, linear, and injective, and its range equals the range of $T$.

\textit{Proof:} \emph{Well-definedness.} If $v+U=w+U$, then $v-w\in U=\ker T$.
Hence
\[
T(v)-T(w)=T(v-w)=0_W,
\]
so $T(v)=T(w)$. Thus the value of $\widetilde{T}$ does not depend on
the representative of the coset.

\emph{Linearity:} For $v,w\in V$ and $a,b\in\mathbb{F}$,
\[
\begin{aligned}
\widetilde{T}\bigl(a(v+U)+b(w+U)\bigr)
&=\widetilde{T}\bigl((av+bw)+U\bigr)\\
&=T(av+bw)\\
&=a\widetilde{T}(v+U)+b\widetilde{T}(w+U).
\end{aligned}
\]

\emph{Injectivity:} If $\widetilde{T}(v+U)=\widetilde{T}(w+U)$,
then $T(v)=T(w)$. Therefore $T(v-w)=0_W$, so $v-w\in U$ and
$v+U=w+U$.

\emph{Range:} Every value $\widetilde{T}(v+U)=T(v)$ belongs to
$\operatorname{range}T$. Conversely, every $y\in\operatorname{range}T$
equals $T(v)=\widetilde{T}(v+U)$ for some $v\in V$. Hence
$\operatorname{range}\widetilde{T}=\operatorname{range}T$.
\hfill$\square$
\end{claim_box}

\begin{claim_box}
\textbf{Factorization:} The quotient map
$\pi:V\longrightarrow V/U$, $\pi(v)=v+U$, satisfies
\[
T=\widetilde{T}\circ\pi.
\]
Moreover, $\widetilde{T}$ is the unique map $F:V/U\longrightarrow W$
with $T=F\circ\pi$: for every coset, this identity forces
$F(v+U)=F(\pi(v))=T(v)$.
\end{claim_box}

\begin{mainbox}{Corollary:}
Restricting the
\emph{codomain} of $\widetilde{T}$ to $\operatorname{range}T$ gives an
isomorphism
\[
V/\ker T\cong\operatorname{range}T.
\]
\end{mainbox}

\begin{claim_box}
\textit{Proof:} Indeed, the induced map is linear and injective by the
proposition,
and is surjective onto this codomain by the equality of ranges.
With its original codomain $W$, $\widetilde{T}$ is surjective (and hence
an isomorphism) if and only if $T$ is surjective. Passing to the quotient
removes the failure of injectivity, but does not enlarge the range.
\hfill$\square$
\end{claim_box}

\begin{claim_box}
\textbf{Alternative proof of rank-nullity (using quotient spaces):}
Suppose now that $V$ is finite-dimensional. Apply the quotient-space
dimension theorem, using $U=\ker T$:
\[
\dim(V/\ker T)+\dim(\ker T)=\dim V.
\]
Since isomorphic spaces have the same dimension,
$\dim(V/\ker T)=\dim(\operatorname{range}T)$. Substituting yields
\[
\dim(\operatorname{range}T)+\dim(\ker T)=\dim V,
\]
which is the rank-nullity theorem. \hfill$\square$
\end{claim_box}

\begin{example_box}
\textbf{Example:} For the earlier map
$T_A:\mathbb{R}^3\longrightarrow\mathbb{R}^2$ given by
$T_A(x,y,z)=(x+2y,-y+z)$, we found
$\ker T_A=\operatorname{span}\{(-2,1,1)\}$. Its nullity is $1$, so
\[
\dim(\operatorname{range}T_A)=3-1=2.
\]
Since its range is a two-dimensional subspace of $\mathbb{R}^2$, the
range is all of $\mathbb{R}^2$ and $T_A$ is surjective.
\end{example_box}

\begin{example_box}
\textbf{Example:} Every linear map from $\mathbb{R}^3$ to
$\mathbb{R}^2$ has a null space of dimension at least $1$, so it sends
some nonzero vector to zero and cannot be injective.
\end{example_box}

\begin{example_box}
\textbf{Example:} Every linear map from $\mathbb{R}^2$ to
$\mathbb{R}^3$ has a range of dimension at most $2$, and therefore cannot
be surjective. For instance, the inclusion
\[
J:\mathbb{R}^2\longrightarrow\mathbb{R}^3,
\qquad J(x,y)=(x,y,0)
\]
has the $xy$-plane as its range and does not attain $(0,0,1)$.
\end{example_box}

\begin{remark_box}
\textbf{Note.} The dimension obstructions above require linearity;
they are not statements about arbitrary maps between the underlying
sets.
\end{remark_box}

\subsection{Matrices}

\begin{definition_box}
\textbf{Definition (matrix):} Let $m,n$ be positive integers. An
\emph{$m\times n$ matrix over $\mathbb{F}$} is a rectangular array of
elements of $\mathbb{F}$ arranged in $m$ rows and $n$ columns:
\[
A=(a_{ij})=
\begin{pmatrix}
a_{11}&a_{12}&\cdots&a_{1n}\\
a_{21}&a_{22}&\cdots&a_{2n}\\
\vdots&\vdots&\ddots&\vdots\\
a_{m1}&a_{m2}&\cdots&a_{mn}
\end{pmatrix}.
\]
The scalar $a_{ij}$ is the \emph{entry} in row $i$ and column $j$.
The size is written with the number of rows first. The set of all such
matrices is denoted by $\mathbb{F}^{m\times n}$ or
$M_{m\times n}(\mathbb{F})$. Two matrices are equal if they have the
same size and all corresponding entries are equal.
\end{definition_box}

\begin{example_box}
\textbf{Example:} The matrix
\[
B=\begin{pmatrix}1&2&0\\3&-1&4\end{pmatrix}
\in\mathbb{R}^{2\times3}
\]
has two rows and three columns. For instance, $b_{21}=3$ and $b_{23}=4$.
Its first row is $(1,2,0)$, and its second column is
$\begin{pmatrix}2\\-1\end{pmatrix}$.
A matrix with $m=n$ is called \emph{square}; a $1\times n$ matrix is a
\emph{row vector}, and an $m\times1$ matrix is a \emph{column vector}.
The \emph{zero matrix} of a given size has every entry equal to zero.
\end{example_box}

\begin{claim_box}
\textbf{Definition (square, diagonal, and triangular matrices):}
A \emph{square matrix of order $n$} has $n$ rows and $n$ columns.
For $A=(a_{ij})\in M_{n\times n}(\mathbb{F})$, its \emph{main diagonal}
consists of $a_{11},a_{22},\ldots,a_{nn}$. We call $A$:
\begin{enumerate}
    \item \emph{Diagonal} if $a_{ij}=0$ whenever $i\neq j$.
    We write $A=\operatorname{diag}(d_1,\ldots,d_n)$ when $a_{ii}=d_i$.
    \item \emph{Upper triangular} if $a_{ij}=0$ whenever $i>j$;
    that is, every entry below the main diagonal is zero.
    \item \emph{Lower triangular} if $a_{ij}=0$ whenever $i<j$;
    that is, every entry above the main diagonal is zero.
    \item A \emph{scalar matrix} if it is diagonal and all its
    diagonal entries are the same scalar $c$. When $c=1$, it is the
    \emph{identity matrix} $I_n$.
\end{enumerate}
Diagonal entries may be zero. A matrix is diagonal if and only if it is
both upper and lower triangular: the two triangular conditions together
say exactly that every off-diagonal entry vanishes. An upper or lower
triangular matrix is called \emph{strictly triangular} of the respective
type if its diagonal entries are also zero.
\end{claim_box}

\begin{example_box}
\textbf{Examples:} Over $\mathbb{R}$, the matrices
\[
D=\begin{pmatrix}2&0&0\\0&0&0\\0&0&-1\end{pmatrix},\qquad
U=\begin{pmatrix}1&2&3\\0&0&4\\0&0&5\end{pmatrix},\qquad
L=\begin{pmatrix}1&0&0\\2&3&0\\4&5&6\end{pmatrix}
\]
are diagonal, upper triangular, and lower triangular, respectively.
The square zero matrix is both diagonal and strictly upper and lower
triangular. In these notes, triangular matrices are square; a rectangular
matrix whose entries satisfy $a_{ij}=0$ for $i>j$ is instead called
\emph{upper trapezoidal}.
\end{example_box}

\begin{claim_box}
\textbf{Definition (matrix of a vector):} Let $V$ be a
finite-dimensional vector space over $\mathbb{F}$ with ordered basis
$\mathcal{B}=(v_1,\ldots,v_n)$. Each $v\in V$ has a unique expansion
$v=x_1 v_1+\cdots+x_n v_n$. The \emph{matrix of $v$ relative to
$\mathcal{B}$}, also called its \emph{coordinate column}, is
\[
{[v]}_{\mathcal{B}}=
\begin{pmatrix}x_1\\\vdots\\x_n\end{pmatrix}
\in\mathbb{F}^{n\times1}.
\]
The order of the entries matches the order of the basis vectors.
We identify this column matrix with a vector of $\mathbb{F}^n$.
It records the coordinates of $v$, not necessarily the vector $v$
itself; changing the basis can change the column.
\end{claim_box}

\begin{example_box}
\textbf{Example (coordinates depend on the basis):} In
$\mathbb{R}^2$, let $v=(3,1)$ and choose
$\mathcal{B}=((1,1),(1,-1))$. Since
\[
(3,1)=2(1,1)+(1,-1),
\]
we have ${[v]}_{\mathcal{B}}=\begin{pmatrix}2\\1\end{pmatrix}$,
whereas its coordinate column in the standard basis is
$\begin{pmatrix}3\\1\end{pmatrix}$.
\end{example_box}

\begin{definition_box}
\textbf{Definition (matrix of a linear map):} Let
$T:V\longrightarrow W$ be linear, where $V,W$ are finite-dimensional
vector spaces over $\mathbb{F}$. Choose ordered bases
\[
\mathcal{B}=(v_1,\ldots,v_n)\quad\text{of }V,
\qquad \mathcal{C}=(w_1,\ldots,w_m)\quad\text{of }W.
\]
For each $j$, express $T(v_j)$ uniquely in the basis $\mathcal{C}$:
\[
T(v_j)=\sum_{i=1}^m a_{ij}w_i.
\]
The \emph{matrix of $T$ relative to $\mathcal{B}$ and $\mathcal{C}$} is
the $m\times n$ matrix
\[
{[T]}_{\mathcal{C}\leftarrow\mathcal{B}}=(a_{ij}).
\]
Its $j$th column is the coordinate column of $T(v_j)$ in the target
basis $\mathcal{C}$. Thus the number of columns is $\dim V$, and the
number of rows is $\dim W$. The arrow in the notation records that
input coordinates use $\mathcal{B}$ and output coordinates use
$\mathcal{C}$.
\end{definition_box}

\begin{mainbox}{Proposition (linear maps as matrix multiplication):}
With the ordered bases above fixed, put
$A={[T]}_{\mathcal{C}\leftarrow\mathcal{B}}$. For every $v\in V$,
\[
{[T(v)]}_{\mathcal{C}}
=A{[v]}_{\mathcal{B}}.
\]
Thus applying $T$ becomes multiplication by $A$ when inputs and outputs
are expressed in their respective bases. Abbreviating these coordinate
matrices by $M(T(v))$, $M(T)$, and $M(v)$ gives
\[
M(T(v))=M(T)M(v).
\]
\end{mainbox}

\begin{claim_box}
\textit{Proof:} Write $v=\sum_{j=1}^n x_j v_j$ and $A=(a_{ij})$.
By linearity and the definition of $A$,
\[
T(v)=\sum_{j=1}^n x_j T(v_j)
=\sum_{i=1}^m\left(\sum_{j=1}^n a_{ij}x_j\right)w_i,
\]
so uniqueness of coordinates gives
\[
{[T(v)]}_{\mathcal{C}}
=\begin{pmatrix}
\sum_{j=1}^n a_{1j}x_j\\
\vdots\\
\sum_{j=1}^n a_{mj}x_j
\end{pmatrix}
=A\begin{pmatrix}x_1\\\vdots\\x_n\end{pmatrix}
=A{[v]}_{\mathcal{B}}.
\]
Here an $m\times n$ matrix multiplies an $n\times1$ column to give
an $m\times1$ column. \hfill$\square$
\end{claim_box}

\begin{example_box}
\textbf{Example (differentiation in polynomial bases):} Consider
$D:\mathcal{P}_2(\mathbb{R})\longrightarrow\mathcal{P}_1(\mathbb{R})$,
$D(p)=p'$, with ordered bases
$\mathcal{B}=(1,x,x^2)$ and $\mathcal{C}=(1,x)$. Since
\[
D(1)=0,\qquad D(x)=1,\qquad D(x^2)=2x,
\]
their coordinate columns give
\[
{[D]}_{\mathcal{C}\leftarrow\mathcal{B}}
=\begin{pmatrix}0&1&0\\0&0&2\end{pmatrix}.
\]
For example, the polynomial $p=a+bx+cx^2$ has coordinate column
${(a,b,c)}^{\mathsf{T}}$, and
\[
\begin{pmatrix}0&1&0\\0&0&2\end{pmatrix}
\begin{pmatrix}a\\b\\c\end{pmatrix}
=\begin{pmatrix}b\\2c\end{pmatrix},
\]
which is the coordinate column of $p'=b+2cx$ in $\mathcal{C}$.
\end{example_box}

\begin{remark_box}
\textbf{Note (the bases matter):} A linear map is not itself a
matrix: a matrix represents it after ordered bases have been chosen.
Changing either basis can change the matrix without changing the map.
For a map $T_A:\mathbb{F}^n\longrightarrow\mathbb{F}^m$, $T_A(x)=Ax$,
the representing matrix in the standard bases is exactly $A$.
\end{remark_box}

\begin{claim_box}
\textbf{Definition (addition and scalar multiplication):} Let
$A=(a_{ij})$ and $B=(b_{ij})$ belong to
$M_{m\times n}(\mathbb{F})$, and let $c\in\mathbb{F}$. Define
\[
A+B=(a_{ij}+b_{ij}),\qquad cA=(ca_{ij}).
\]
Thus addition adds corresponding entries, and scalar multiplication
multiplies every entry by the same scalar. Both operations preserve the
size of the matrix; addition requires the matrices to have the same size.
\end{claim_box}

\begin{example_box}
\textbf{Example:} Over $\mathbb{R}$,
\[
\begin{pmatrix}1&2\\0&-1\end{pmatrix}
+\begin{pmatrix}3&0\\4&2\end{pmatrix}
=\begin{pmatrix}4&2\\4&1\end{pmatrix},
\qquad
3\begin{pmatrix}1&2\\0&-1\end{pmatrix}
=\begin{pmatrix}3&6\\0&-3\end{pmatrix}.
\]
Under these operations, $M_{m\times n}(\mathbb{F})$ is a vector space:
the axioms hold entrywise because they hold in $\mathbb{F}$. Its zero
vector is the zero matrix, and the additive inverse of $A$ is $-A$.
\end{example_box}

\begin{claim_box}
\textbf{Proposition (matrix representation respects linear combinations):}
Fix ordered bases $\mathcal{B}=(v_1,\ldots,v_n)$ of $V$ and
$\mathcal{C}=(w_1,\ldots,w_m)$ of $W$. For
$T\in\mathcal{L}(V,W)$, abbreviate
\[
M(T)={[T]}_{\mathcal{C}\leftarrow\mathcal{B}}.
\]
Then for all $S,T\in\mathcal{L}(V,W)$ and $c\in\mathbb{F}$,
\[
M(S+T)=M(S)+M(T),\qquad M(cT)=cM(T).
\]

\textit{Proof:} Write $M(S)=(s_{ij})$ and $M(T)=(t_{ij})$. By the
definition of a representing matrix, for each $j$,
\[
S(v_j)=\sum_{i=1}^m s_{ij}w_i,
\qquad T(v_j)=\sum_{i=1}^m t_{ij}w_i.
\]
Using pointwise addition of maps gives
\[
(S+T)(v_j)=S(v_j)+T(v_j)
=\sum_{i=1}^m(s_{ij}+t_{ij})w_i.
\]
Uniqueness of coordinates in $\mathcal{C}$ shows that the entry in row
$i$, column $j$ of $M(S+T)$ is $s_{ij}+t_{ij}$. These are exactly the
entries of $M(S)+M(T)$, proving the first identity.

Similarly, pointwise scalar multiplication gives
\[
(cT)(v_j)=cT(v_j)=\sum_{i=1}^m(ct_{ij})w_i.
\]
Thus the entry in row $i$, column $j$ of $M(cT)$ is $ct_{ij}$, which
is the corresponding entry of $cM(T)$. This proves the second identity.
\hfill$\square$
\end{claim_box}

\begin{claim_box}
\textbf{Consequence:} With these bases fixed, the assignment
\[
M:\mathcal{L}(V,W)\longrightarrow M_{m\times n}(\mathbb{F}),
\qquad T\longmapsto M(T),
\]
is itself a linear map. In particular,
$M(aS+bT)=aM(S)+bM(T)$ for all $a,b\in\mathbb{F}$.
All matrices in these identities must use the same ordered domain and
codomain bases.
\end{claim_box}

\begin{definition_box}
\textbf{Definition (matrix multiplication):} Let
$A=(a_{ik})\in M_{m\times n}(\mathbb{F})$ and
$B=(b_{kj})\in M_{n\times p}(\mathbb{F})$. Their \emph{product}
$AB$ is the $m\times p$ matrix whose entries are
\[
{(AB)}_{ij}=\sum_{k=1}^n a_{ik}b_{kj},
\qquad 1\leq i\leq m,\quad 1\leq j\leq p.
\]
To find the entry in row $i$, column $j$, multiply corresponding entries
of row $i$ of $A$ and column $j$ of $B$, then add the products.
The number of columns of $A$ must equal the number of rows of $B$:
\[
(m\times n)(n\times p)\longrightarrow m\times p.
\]
This is not entrywise multiplication.
\end{definition_box}

\begin{example_box}
\textbf{Example (row-by-column multiplication):} Over $\mathbb{R}$,
let
\[
A=\begin{pmatrix}1&2&0\\0&-1&1\end{pmatrix},
\qquad B=\begin{pmatrix}1&0\\2&1\\3&4\end{pmatrix}.
\]
Since $A$ is $2\times3$ and $B$ is $3\times2$, the product $AB$ is
$2\times2$. Computing each entry gives
\[
\begin{aligned}
AB
&=\begin{pmatrix}
1\cdot1+2\cdot2+0\cdot3 & 1\cdot0+2\cdot1+0\cdot4\\
0\cdot1+(-1)\cdot2+1\cdot3 & 0\cdot0+(-1)\cdot1+1\cdot4
\end{pmatrix}\\[4pt]
&=\begin{pmatrix}5&2\\1&3\end{pmatrix}.
\end{aligned}
\]
\end{example_box}

\begin{remark_box}
\textbf{Note:} The reverse product $BA$ need not
be defined. Even when both products exist, they need not be equal;
in this example $BA$ is $3\times3$, whereas $AB$ is $2\times2$.
\end{remark_box}

\begin{claim_box}
\textbf{Definition (identity matrix):} The \emph{identity matrix}
$I_n\in\mathbb{F}^{n\times n}$ has $1$ in every diagonal position
and $0$ elsewhere. Write $0_{m\times n}$ for the zero matrix of size
$m\times n$.
\end{claim_box}

\begin{claim_box}
\textbf{Properties of matrix addition and scalar multiplication:}
For $A,B,C\in\mathbb{F}^{m\times n}$ and $a,b\in\mathbb{F}$,
\[
\begin{aligned}
A+B&=B+A, & (A+B)+C&=A+(B+C),\\
A+0_{m\times n}&=A, & A+(-A)&=0_{m\times n},\\
a(A+B)&=aA+aB, & (a+b)A&=aA+bA,\\
a(bA)&=(ab)A, & 1A&=A,\qquad 0A=0_{m\times n}.
\end{aligned}
\]
These identities hold entry by entry because the corresponding
identities hold in the field $\mathbb{F}$.
\end{claim_box}

\begin{claim_box}
\textbf{Properties of matrix multiplication:} All products and
sums below require the indicated compatible sizes.
\begin{enumerate}
    \item \emph{Associativity.} If $A$ is $m\times n$, $B$ is
    $n\times p$, and $C$ is $p\times q$, then
    \[
    (AB)C=A(BC)\in\mathbb{F}^{m\times q}.
    \]
    Thus parentheses may be omitted from a product $ABC$, while
    the order of the factors must be retained.
    \item \emph{Distributivity.} If $A$ is $m\times n$ and $B,C$
    are $n\times p$, then
    \[
    A(B+C)=AB+AC.
    \]
    If $A,B$ are $m\times n$ and $C$ is $n\times p$, then
    \[
    (A+B)C=AC+BC.
    \]
    \item \emph{Compatibility with scalars.} If $A$ is $m\times n$,
    $B$ is $n\times p$, and $c\in\mathbb{F}$, then
    \[
    c(AB)=(cA)B=A(cB).
    \]
    \item \emph{Identity and zero.} For $A\in\mathbb{F}^{m\times n}$,
    \[
    I_m A=A=A I_n,\qquad
    A0_{n\times p}=0_{m\times p},\qquad
    0_{q\times m}A=0_{q\times n}.
    \]
\end{enumerate}

\textit{Proof:} For associativity, the entry in row $i$, column $j$ is
\[
\begin{aligned}
{((AB)C)}_{ij}
&=\sum_{k=1}^p\left(\sum_{\ell=1}^n a_{i\ell}b_{\ell k}\right)c_{kj}\\
&=\sum_{\ell=1}^n a_{i\ell}
  \left(\sum_{k=1}^p b_{\ell k}c_{kj}\right)
={\bigl(A(BC)\bigr)}_{ij}.
\end{aligned}
\]
Distributivity follows by expanding each entry, for example
\[
{(A(B+C))}_{ij}
=\sum_{k=1}^n a_{ik}(b_{kj}+c_{kj})
={(AB)}_{ij}+{(AC)}_{ij}.
\]
Expanding $(a_{ik}+b_{ik})c_{kj}$ proves the other distributive law.
A scalar factor can be moved outside each finite sum, proving scalar
compatibility. Finally, multiplication by an identity matrix selects
the corresponding row or column of $A$, while multiplication by a zero
matrix gives only zero summands. These entrywise equalities prove all
the identities. \hfill$\square$
\end{claim_box}

\begin{claim_box}
\textbf{Definition (left and right inverses):} Let
$A\in\mathbb{F}^{m\times n}$. A matrix $L\in\mathbb{F}^{n\times m}$
is a \emph{left inverse} of $A$ if $LA=I_n$. A matrix
$R\in\mathbb{F}^{n\times m}$ is a \emph{right inverse} of $A$ if
$AR=I_m$. The side refers to the position of the inverse in the product.
Both kinds of inverse necessarily have size $n\times m$; their having
the same size alone does not imply their existence or equality.
\end{claim_box}

\begin{definition_box}
\textbf{Definition (inverse and invertibility):} A square matrix
$A\in\mathbb{F}^{n\times n}$ is \emph{invertible}, or
\emph{nonsingular}, if there is a matrix $B\in\mathbb{F}^{n\times n}$
such that
\[
AB=BA=I_n.
\]
Such a $B$ is a \emph{two-sided inverse} of $A$. We will prove that it
is unique and denote it by $A^{-1}$. A square matrix with no inverse is
called \emph{singular}. Matrix inversion is not entrywise reciprocation.
\end{definition_box}

\begin{claim_box}
\textbf{Theorem (left and right inverses agree):} If a matrix
$A\in\mathbb{F}^{m\times n}$ admits both a left inverse $L$ and a
right inverse $R$, then $L=R$. Moreover, $m=n$, so $A$ is invertible
and this common inverse is unique.

\textit{Proof:} Associativity and the identity laws give
\[
L=L I_m=L(AR)=(LA)R=I_n R=R.
\]
To see why $A$ must be square, consider
$T_A:\mathbb{F}^n\longrightarrow\mathbb{F}^m$, $T_A(x)=Ax$.
If $Ax=0$, then $x=LAx=0$, so $T_A$ is injective and $n\leq m$.
For every $y\in\mathbb{F}^m$, we have $y=A(Ry)$, so $T_A$ is
surjective and $m\leq n$. Thus $m=n$.
Finally, any other left inverse $L'$ equals $R$ by the same calculation,
and any other right inverse $R'$ equals $L$. In particular, any two
two-sided inverses coincide. \hfill$\square$
\end{claim_box}

\begin{mainbox}{Theorem}
Let
$A\in\mathbb{F}^{n\times n}$. Then
\[
A\text{ has a left inverse}
\quad\Longleftrightarrow\quad
A\text{ has a right inverse}.
\]
Whenever these equivalent conditions hold, $A$ is invertible and
every one-sided inverse equals the unique two-sided inverse $A^{-1}$.
\end{mainbox}

\begin{claim_box}
\textit{Proof:} $(\Rightarrow)$ Suppose $L$ is a left inverse,
so $LA=I_n$.
If $Ax=0$, then
\[
x=I_n x=(LA)x=L(Ax)=0.
\]
Thus $T_A:\mathbb{F}^n\longrightarrow\mathbb{F}^n$, $T_A(x)=Ax$,
is injective. Since the domain and codomain have the same finite
dimension, rank-nullity implies that $T_A$ is also surjective.
For any $y\in\mathbb{F}^n$, choose $x$ with $Ax=y$. Then
\[
ALy=AL(Ax)=A(LA)x=Ax=y.
\]
Hence $AL$ acts as the identity on every vector, and evaluating on
the standard basis gives $AL=I_n$. Therefore the very same matrix $L$
is also a right inverse, and $A^{-1}=L$.

$(\Leftarrow)$ Suppose $R$ is a right inverse, so $AR=I_n$. Every
$y\in\mathbb{F}^n$ satisfies $y=A(Ry)$, making $T_A$ surjective
and hence injective by equal finite dimensions. For every $x$,
\[
A(RAx-x)=(AR)Ax-Ax=0.
\]
Injectivity implies $RAx-x=0$. Thus $RA=I_n$, so $R$ is also a
left inverse.

This proves both implications. In either case the given inverse is
two-sided, and uniqueness follows from the preceding theorem.
\hfill$\square$
\end{claim_box}

\begin{remark_box}
\textbf{Caution (one-sided inverses need not be unique):}
For rectangular matrices, a left inverse alone or a right inverse alone
need not be unique. For example, for every $t\in\mathbb{F}$,
\[
\begin{pmatrix}1&t\end{pmatrix}\begin{pmatrix}1\\0\end{pmatrix}
=I_1,\qquad
\begin{pmatrix}1&0\end{pmatrix}\begin{pmatrix}1\\t\end{pmatrix}
=I_1.
\]
Thus the column matrix on the left has distinct left inverses, and
the row matrix on the right has distinct right inverses (take $t=0,1$).
Neither rectangular matrix can have both kinds of inverse.
\end{remark_box}

\begin{example_box}
\textbf{Example (a zero product with nonzero factors):} Unlike
scalars in a field, nonzero matrices can multiply to zero. For example,
\[
A=\begin{pmatrix}1&0\\0&0\end{pmatrix},\qquad
B=\begin{pmatrix}0&0\\1&0\end{pmatrix}
\quad\Longrightarrow\quad
AB=\begin{pmatrix}0&0\\0&0\end{pmatrix}.
\]
Both $A$ and $B$ are nonzero. Moreover, $BA=B\neq0$, so $AB=0$
does not imply $BA=0$, even when both products are defined.
\end{example_box}

\begin{claim_box}
\textbf{Theorem:} Let
$A\in\mathbb{F}^{m\times n}$ and $B\in\mathbb{F}^{n\times p}$.
Write $T_A(x)=Ax$ and $T_B(y)=By$. Then
\[
AB=0_{m\times p}
\quad\Longleftrightarrow\quad
\operatorname{range}T_B\subseteq\ker T_A.
\]
Equivalently, every column of $B$ lies in the null space of $A$.
Consequently, if $AB=0$, then
\[
\operatorname{rank}A+\operatorname{rank}B\leq n.
\]

\textit{Proof:} If $AB=0$, every vector $z=By$ in the range of $T_B$
satisfies $Az=ABy=0$, so $z\in\ker T_A$. Conversely, the inclusion
implies $ABy=0$ for every $y\in\mathbb{F}^p$. Taking $y$ to be each
standard basis vector shows that every column of $AB$ is zero.
The column formulation follows because column $j$ of $AB$ is $AB_j$,
where $B_j$ is column $j$ of $B$. Finally, the inclusion and
rank-nullity give
\[
\operatorname{rank}B=\dim(\operatorname{range}T_B)
\leq\dim(\ker T_A)=n-\operatorname{rank}A.
\]
Rearranging proves the inequality. \hfill$\square$
\end{claim_box}

\begin{claim_box}
\textbf{Proposition:} For the same
matrix sizes, suppose $AB=0$.
\begin{enumerate}
    \item If $A$ has a left inverse, then $B=0$.
    In particular, this holds if $A$ is square and invertible.
    \item If $B$ has a right inverse, then $A=0$.
    In particular, this holds if $B$ is square and invertible.
\end{enumerate}

\textit{Proof:} If $LA=I_n$, then
\[
B=I_n B=(LA)B=L(AB)=0.
\]
If $BR=I_n$, then
\[
A=A I_n=A(BR)=(AB)R=0.
\]
Thus cancellation is justified by an inverse on the appropriate side,
not merely by a factor being nonzero. \hfill$\square$
\end{claim_box}

\begin{claim_box}
\textbf{Corollary:} If $A$ has a left inverse
and $AB=AC$, where $B,C$ have the same compatible size, then $B=C$.
If $A$ has a right inverse and $BA=CA$, then $B=C$.

\textit{Proof:} Distributivity gives $A(B-C)=0$ in the first case
and $(B-C)A=0$ in the second. Apply the preceding proposition to
obtain $B-C=0$. \hfill$\square$
\end{claim_box}

\begin{claim_box}
\textbf{Proposition:}
Let $A\in\mathbb{F}^{n\times n}$ with $n\geq1$. Then $A$ is singular
if and only if there exists a nonzero $B\in\mathbb{F}^{n\times n}$
such that $AB=0$.

\textit{Proof:} If $A$ is singular, $T_A$ cannot be injective:
an injective operator on $\mathbb{F}^n$ is bijective, and its linear
inverse would provide an inverse matrix for $A$. Choose a nonzero
$v\in\ker T_A$ and let $B$ have first column $v$ and all other
columns zero. Then $B\neq0$ and every column of $AB$ is zero.
Conversely, if $AB=0$ with $B\neq0$, $A$ cannot be invertible by
the zero-factor proposition. Hence $A$ is singular. \hfill$\square$
\end{claim_box}

\begin{claim_box}
\textbf{Consequence:} If $A,B$ are square matrices of the same
order, $AB=0$, and both factors are nonzero, then both factors are
singular. Indeed, invertibility of either would force the other to be
zero. The converse is false: two singular matrices need not have zero
product. For instance,
$P=\begin{pmatrix}1&0\\0&0\end{pmatrix}$ is singular but $P^2=P\neq0$.
\end{claim_box}

\begin{mainbox}{Theorem (inverse of a $2\times2$ matrix)}
Let
\[
A=\begin{pmatrix}a&b\\c&d\end{pmatrix},\qquad a,b,c,d\in\mathbb{F}.
\]
The scalar $\det A=ad-bc$ is called the \emph{determinant} of this
$2\times2$ matrix. Then $A$ is invertible if and only if $ad-bc\neq0$,
and in that case
\[
\boxed{A^{-1}=\frac{1}{ad-bc}\begin{pmatrix}d&-b\\-c&a\end{pmatrix}.}
\]
To compute it, interchange the diagonal entries, negate the two
off-diagonal entries, and multiply by ${(ad-bc)}^{-1}$.
\end{mainbox}

\begin{claim_box}
\textit{Proof:} Put $C=\begin{pmatrix}d&-b\\-c&a\end{pmatrix}$.
Direct multiplication gives
\[
AC=\begin{pmatrix}ad-bc&0\\0&ad-bc\end{pmatrix}
=CA=(ad-bc)I_2.
\]
If $ad-bc\neq0$, division by this scalar gives a two-sided inverse,
proving the formula.

If $ad-bc=0$, then $AC=0$. Were $A$ invertible, multiplying by
$A^{-1}$ on the left would give $C=0$. The entries of $C$ then force
$a=b=c=d=0$, so $A=0$, which cannot have an inverse because $0B=0\neq I_2$
for every $B$. This contradiction proves the converse. \hfill$\square$
\end{claim_box}

\begin{example_box}
\textbf{Example (computing a $2\times2$ inverse):} Over
$\mathbb{R}$, let
\[
A=\begin{pmatrix}2&1\\5&3\end{pmatrix}.
\]
Since $\det A=2\cdot3-1\cdot5=1$, the formula gives
\[
A^{-1}=\begin{pmatrix}3&-1\\-5&2\end{pmatrix}.
\]
Indeed,
\[
\begin{pmatrix}2&1\\5&3\end{pmatrix}
\begin{pmatrix}3&-1\\-5&2\end{pmatrix}
=\begin{pmatrix}1&0\\0&1\end{pmatrix}.
\]
In contrast, $\begin{pmatrix}1&2\\2&4\end{pmatrix}$ has determinant
zero and hence no inverse.
\end{example_box}

\begin{problem_box}
\textbf{Exercise:}
Let $n\geq1$ be an integer, and define $A=(a_{ij})\in\mathbb{R}^{n\times n}$
by
\[
a_{ij}=\frac{1}{i+j-1},\qquad 1\leq i,j\leq n.
\]
Prove that:
\begin{enumerate}
    \item $A$ is invertible for every $n\geq1$.
    \item Every entry of $A^{-1}$ is an integer; that is,
    $A^{-1}\in M_{n\times n}(\mathbb{Z})$.
\end{enumerate}
\end{problem_box}

\begin{claim_box}
\textbf{Solution:} We construct a matrix
$B$ with integer entries such that $AB=I_n$. Since $A$ is square,
this will prove both invertibility and $A^{-1}=B$. The only polynomial
fact we need is that a polynomial of degree at most $n-1$ with $n$
distinct roots is identically zero. Empty products below equal $1$,
and $0!=1$.

\emph{Construct one column of the inverse:} Fix
$j\in\{1,\ldots,n\}$. We want numbers $b_{1j},\ldots,b_{nj}$ such that
\[
\sum_{i=1}^n\frac{b_{ij}}{k+i-1}
=\begin{cases}1,&k=j,\\0,&k\neq j,\end{cases}
\qquad 1\leq k\leq n.
\]
These are exactly the equations for column $j$ of $AB=I_n$.
Introduce the polynomials
\[
Q(x)=\prod_{r=1}^n(x+r-1),\qquad
P_j(x)=Q(j)\prod_{\substack{1\leq k\leq n\\k\neq j}}
\frac{x-k}{j-k}.
\]
The polynomial $P_j$ has degree $n-1$, and its factors give
$P_j(j)=Q(j)$ and $P_j(k)=0$ for all other $k\in\{1,\ldots,n\}$.
Thus $P_j(k)/Q(k)$ has exactly the desired values.

\emph{Write this quotient as a sum:} Set
\[
b_{ij}=\frac{P_j(1-i)}{
\displaystyle\prod_{\substack{1\leq r\leq n\\r\neq i}}(r-i)}.
\]
Every denominator here is nonzero. We claim the polynomial identity
\[
P_j(x)=\sum_{i=1}^n b_{ij}
\prod_{\substack{1\leq r\leq n\\r\neq i}}(x+r-1).
\]
To check it, substitute $x=1-s$, where $1\leq s\leq n$.
All terms on the right vanish except $i=s$, and that term equals
$P_j(1-s)$ by the definition of $b_{sj}$. Both sides have degree
at most $n-1$, so their difference has $n$ distinct roots and is zero.
Dividing by $Q(x)$, whenever $Q(x)\neq0$, gives
\[
\frac{P_j(x)}{Q(x)}=\sum_{i=1}^n\frac{b_{ij}}{x+i-1}.
\]
In particular, substitution of $x=k\in\{1,\ldots,n\}$ proves the
equations in Step 1. Doing this for every $j$ constructs $B$ with
$AB=I_n$, so $A$ is invertible and $B=A^{-1}$.

\emph{Show that the entries are integers:} The products above
can be written using factorials:
\[
\begin{aligned}
Q(j)&=\frac{(n+j-1)!}{(j-1)!},\\
\prod_{\substack{1\leq k\leq n\\k\neq j}}(j-k)
&={(-1)}^{n-j}(j-1)!(n-j)!,\\
\prod_{\substack{1\leq k\leq n\\k\neq j}}(1-i-k)
&={(-1)}^{n-1}\frac{(n+i-1)!}{(i-1)!(i+j-1)},\\
\prod_{\substack{1\leq r\leq n\\r\neq i}}(r-i)
&={(-1)}^{i-1}(i-1)!(n-i)!.
\end{aligned}
\]
Substitution into the formula for $b_{ij}$ yields
\[
b_{ij}={(-1)}^{i+j}
\frac{(n+i-1)!(n+j-1)!}
{(i+j-1){((i-1)!)}^2{((j-1)!)}^2(n-i)!(n-j)!}.
\]
Although this expression is a quotient, it can be regrouped as
\[
\boxed{
b_{ij}={(-1)}^{i+j}(i+j-1)
\binom{n+i-1}{n-j}\binom{n+j-1}{n-i}
\binom{i+j-2}{i-1}^{\!2}.}
\]
Indeed, expanding the three binomial coefficients into factorials
recovers the preceding expression. All their upper and lower arguments
are nonnegative integers with the lower argument no larger than the
upper one. Each binomial coefficient is therefore an integer, so
$b_{ij}\in\mathbb{Z}$. Since $B=A^{-1}$, this completes the proof.
\hfill$\square$
\end{claim_box}

\begin{mainbox}{Theorem (the matrix of a composition)}
Let
$T:U\longrightarrow V$ and $S:V\longrightarrow W$ be linear maps
between finite-dimensional vector spaces over $\mathbb{F}$. Fix ordered
bases
\[
\mathcal{B}=(u_1,\ldots,u_p),\qquad
\mathcal{C}=(v_1,\ldots,v_n),\qquad
\mathcal{D}=(w_1,\ldots,w_m)
\]
of $U,V,W$, respectively. Then
\[
{[ST]}_{\mathcal{D}\leftarrow\mathcal{B}}
={[S]}_{\mathcal{D}\leftarrow\mathcal{C}}
{[T]}_{\mathcal{C}\leftarrow\mathcal{B}}.
\]
Abbreviating these matrices by $M(ST),M(S),M(T)$, this says
\[
M(ST)=M(S)M(T).
\]
Here $M(S)$ is $m\times n$, $M(T)$ is $n\times p$, and their product
is $m\times p$, as required for the matrix of $ST:U\longrightarrow W$.
\end{mainbox}

\begin{claim_box}
\textit{Proof:} Write $M(S)=(a_{ik})$ and $M(T)=(b_{kj})$.
By the definition of these matrices,
\[
T(u_j)=\sum_{k=1}^n b_{kj}v_k,
\qquad S(v_k)=\sum_{i=1}^m a_{ik}w_i.
\]
For each domain basis vector $u_j$, linearity of $S$ gives
\[
\begin{aligned}
(ST)(u_j)
&=S\left(\sum_{k=1}^n b_{kj}v_k\right)\\
&=\sum_{k=1}^n b_{kj}S(v_k)\\
&=\sum_{k=1}^n\sum_{i=1}^m b_{kj}a_{ik}w_i\\
&=\sum_{i=1}^m\left(\sum_{k=1}^n a_{ik}b_{kj}\right)w_i.
\end{aligned}
\]
Thus the entry in row $i$, column $j$ of $M(ST)$ is
$\sum_{k=1}^n a_{ik}b_{kj}$. By the definition of matrix multiplication,
this is precisely the corresponding entry of $M(S)M(T)$.
Since all corresponding entries agree, $M(ST)=M(S)M(T)$.
\hfill$\square$
\end{claim_box}

\begin{remark_box}
\textbf{Note:} The same intermediate basis $\mathcal{C}$ must be
used for the output of $T$ and the input of $S$. The order also matches
composition: $T$ acts first, so its matrix appears on the right.
\end{remark_box}

\begin{definition_box}
\textbf{Definition (row echelon form and RREF):} The
\emph{leading entry} of a nonzero row is its leftmost nonzero entry.
A matrix is in \emph{row echelon form} (REF) if:
\begin{enumerate}
    \item All zero rows occur below all nonzero rows.
    \item The leading entry of each nonzero row is strictly to the
    right of the leading entry of the preceding nonzero row.
    \item Every entry below a leading entry is zero.
\end{enumerate}
It is in \emph{reduced row echelon form} (RREF) if it also satisfies:
\begin{enumerate}
    \item Every leading entry is $1$.
    \item Each leading $1$ is the only nonzero entry in its column.
\end{enumerate}
Leading entries need not lie on the main diagonal: columns may be
skipped. The zero matrix satisfies these conditions as well.
\end{definition_box}

\begin{claim_box}
\textbf{Proposition (echelon form and triangularity).} Every
square matrix in REF is upper triangular. More generally, every
rectangular matrix in REF is upper trapezoidal.

\textit{Proof:} Let $A$ have $r$ nonzero rows. They are the first $r$
rows because all zero rows are at the bottom. For $1\leq i\leq r$,
let $p_i$ be the column index of the leading entry in row $i$.
The echelon conditions give
\[
1\leq p_1<p_2<\cdots<p_r,
\]
so $p_i\geq i$. If $j<i$, then $j<p_i$, and hence $a_{ij}=0$
because all entries before the leading entry of row $i$ vanish.
For $i>r$, the entire row is zero. Thus $a_{ij}=0$ whenever $i>j$,
which proves both assertions. \hfill$\square$
\end{claim_box}

\begin{claim_box}
\textbf{Theorem (square RREF with no zero rows):} If
$A\in M_{n\times n}(\mathbb{F})$ is in RREF and has no zero row,
then $A=I_n$.

\textit{Proof:} Every row has a leading entry, so there are $n$
leading columns with indices
\[
1\leq p_1<p_2<\cdots<p_n\leq n.
\]
The only possibility is $p_i=i$ for every $i$. Thus every diagonal
entry is a leading $1$. Since $A$ is reduced, that $1$ is the only
nonzero entry in its column. Consequently, all off-diagonal entries
are zero, and $A=I_n$. \hfill$\square$
\end{claim_box}

\begin{remark_box}
\textbf{Caution (the hypotheses matter):} An upper triangular
matrix need not be in REF\@: for example,
$\begin{pmatrix}0&1\\0&1\end{pmatrix}$ has its two leading entries
in the same column. Also, the square assumption in the preceding
theorem is essential. The rectangular matrix
\[
\begin{pmatrix}1&0&2\\0&1&3\end{pmatrix}
\]
is in RREF and has no zero row, but is not an identity matrix.
\end{remark_box}

\begin{definition_box}
\textbf{Definition (transpose of a matrix):} The
\emph{transpose} of $A=(a_{ij})\in\mathbb{F}^{m\times n}$ is the
$n\times m$ matrix $A^{\mathsf{T}}$ defined by
\[
{(A^{\mathsf{T}})}_{ji}=a_{ij},
\qquad 1\leq i\leq m,\quad 1\leq j\leq n.
\]
It interchanges rows and columns without conjugating entries, even
over $\mathbb{C}$. The notation $A'$ is also used for the transpose;
here we use $A^{\mathsf{T}}$ to make its meaning explicit. For example,
\[
A=\begin{pmatrix}1&2&3\\4&5&6\end{pmatrix}
\quad\Longrightarrow\quad
A^{\mathsf{T}}=\begin{pmatrix}1&4\\2&5\\3&6\end{pmatrix}.
\]
Transposing twice restores the original matrix, and transposing an
identity matrix leaves it unchanged:
\[
{(A^{\mathsf{T}})}^{\mathsf{T}}=A,\qquad I_n^{\mathsf{T}}=I_n.
\]
Both identities follow directly by interchanging the row and column
indices in the definition.
\end{definition_box}

\begin{claim_box}
\textbf{Proposition (transpose of a product):} If $A$ is
$m\times n$ and $B$ is $n\times p$, then
\[
{(AB)}^{\mathsf{T}}=B^{\mathsf{T}}A^{\mathsf{T}}.
\]
In particular, transposition reverses the order of multiplication.

\textit{Proof.} Both sides are $p\times m$. For each $i,j$,
\[
\begin{aligned}
{\bigl({(AB)}^{\mathsf{T}}\bigr)}_{ij}
&={(AB)}_{ji}=\sum_{k=1}^n a_{jk}b_{ki}\\
&=\sum_{k=1}^n {(B^{\mathsf{T}})}_{ik}{(A^{\mathsf{T}})}_{kj}
={(B^{\mathsf{T}}A^{\mathsf{T}})}_{ij}.
\end{aligned}
\]
Equality of all entries proves the formula. \hfill$\square$
\end{claim_box}

\begin{claim_box}
\textbf{Theorem (transpose and inverse):} If
$A\in\mathbb{F}^{n\times n}$ is invertible, then $A^{\mathsf{T}}$
is invertible and
\[
\boxed{{(A^{\mathsf{T}})}^{-1}={(A^{-1})}^{\mathsf{T}}.}
\]
In prime notation, this is ${(A')}^{-1}={(A^{-1})}'$.

\textit{Proof.} Transpose the identities $A^{-1}A=I_n$ and
$AA^{-1}=I_n$, using the product rule:
\[
A^{\mathsf{T}}{(A^{-1})}^{\mathsf{T}}=I_n,\qquad
{(A^{-1})}^{\mathsf{T}}A^{\mathsf{T}}=I_n.
\]
Thus ${(A^{-1})}^{\mathsf{T}}$ is a two-sided inverse of
$A^{\mathsf{T}}$. Uniqueness of the inverse proves the result.
\hfill$\square$
\end{claim_box}

\begin{claim_box}
\textbf{Definition (column echelon form and RCEF):}
A matrix $A$ is in \emph{column echelon
form} if $A^{\mathsf{T}}$ is in row echelon form, and in
\emph{reduced column echelon form} (RCEF) if $A^{\mathsf{T}}$ is in RREF\@.
Explicitly, RCEF means that zero columns are on the right, the topmost
nonzero entry of each nonzero column is $1$, these leading $1$s move
strictly downward as the columns move right, and each is the only
nonzero entry in its row.
\end{claim_box}

\begin{example_box}
\textbf{Example (the two reduced forms):} Over $\mathbb{R}$,
\[
R=\begin{pmatrix}1&2&0&3\\0&0&1&-1\\0&0&0&0\end{pmatrix}
\quad\text{is in RREF},\qquad
R^{\mathsf{T}}=
\begin{pmatrix}1&0&0\\2&0&0\\0&1&0\\3&-1&0\end{pmatrix}
\quad\text{is in RCEF}.
\]
The leading entries of $R$ are in positions $(1,1)$ and $(2,3)$;
those of $R^{\mathsf{T}}$ are in positions $(1,1)$ and $(3,2)$.
RREF does not require every column to have a leading $1$, nor does
RCEF require every row to have one.
\end{example_box}

\begin{remark_box}
\textbf{Note (row reduction versus column reduction):}
Elementary row operations produce RREF\@; the corresponding operations
on columns produce RCEF\@. Equivalently, reduce $A^{\mathsf{T}}$ by rows
and transpose the result to obtain the RCEF of $A$. These are different
operations: to solve $Ax=b$ in the original variables, use row
operations on the entire augmented matrix $[A\mid b]$. Column
operations on $A$ generally change the variables and require a
corresponding change of coordinates; they cannot simply replace row
operations on an augmented system.
\end{remark_box}

\begin{definition_box}
\textbf{Definition (elementary matrices):} An \emph{elementary
matrix of order $m$} is a matrix obtained from $I_m$ by one elementary
row operation. There are three types:
\begin{enumerate}
    \item Interchange rows $i$ and $j$, where $i\neq j$.
    \item Multiply row $i$ by a nonzero scalar $c$.
    \item Add $c$ times row $j$ to row $i$, where $i\neq j$.
\end{enumerate}
In the second type, the matrix differs from $I_m$ only in its
$(i,i)$ entry, which is $c$. In the third type, it differs only in
its $(i,j)$ entry, which is $c$ instead of zero. The order of the
elementary matrix matches the number of rows of the matrix on which
the operation will act.
\end{definition_box}

\begin{example_box}
\textbf{Examples:} The elementary matrices
\[
P=\begin{pmatrix}0&1\\1&0\end{pmatrix},\qquad
D=\begin{pmatrix}c&0\\0&1\end{pmatrix}\ (c\neq0),\qquad
H=\begin{pmatrix}1&0\\c&1\end{pmatrix}
\]
are obtained from $I_2$ by $R_1\leftrightarrow R_2$,
$R_1\leftarrow cR_1$, and $R_2\leftarrow R_2+cR_1$, respectively.
For an arbitrary $2\times n$ matrix with rows $u$ and $v$,
\[
P\begin{pmatrix}u\\v\end{pmatrix}=\begin{pmatrix}v\\u\end{pmatrix},
\qquad
D\begin{pmatrix}u\\v\end{pmatrix}=\begin{pmatrix}cu\\v\end{pmatrix},
\qquad
H\begin{pmatrix}u\\v\end{pmatrix}=\begin{pmatrix}u\\v+cu\end{pmatrix}.
\]
\end{example_box}

\begin{claim_box}
\textbf{Theorem (row operations are left multiplication):}
Let $A\in\mathbb{F}^{m\times n}$, and let $E$ be obtained from $I_m$
by an elementary row operation. Then $EA$ is exactly the matrix
obtained by performing that same operation on $A$.

\textit{Proof:} Write $A_{k,*}$ for row $k$ of $A$. Matrix
multiplication says that row $i$ of $EA$ is
\[
{(EA)}_{i,*}=\sum_{k=1}^m e_{ik}A_{k,*}.
\]
Thus each row of $E$ specifies a linear combination of the rows of $A$.
For an unchanged row of $I_m$, this formula returns the corresponding
row of $A$. It remains to check the changed rows:
\begin{enumerate}
    \item If rows $i,j$ of $I_m$ are swapped, row $i$ of $E$ selects
    $A_{j,*}$ and row $j$ selects $A_{i,*}$. Hence $EA$ swaps those rows.
    \item If row $i$ is multiplied by $c$, its only nonzero entry
    is $e_{ii}=c$. The formula gives $cA_{i,*}$.
    \item If $c$ times row $j$ is added to row $i$, then row $i$
    of $E$ has entries $e_{ii}=1$ and $e_{ij}=c$, with all others
    zero. The formula gives $A_{i,*}+cA_{j,*}$.
\end{enumerate}
These are precisely the three stated operations. \hfill$\square$
\end{claim_box}

\begin{claim_box}
\textbf{Corollary (elementary matrices are invertible):}
The inverse of an elementary matrix is the elementary matrix for the
reverse operation: swap the same rows again, scale by $c^{-1}$,
or add $-c$ times the other row, respectively.

\textit{Proof.} Let $E$ represent an operation and $F$ its reverse.
By the theorem, $FE$ applies the reverse operation to $E$, giving
$I_m$. Similarly, $EF=I_m$. Thus $F=E^{-1}$. \hfill$\square$
\end{claim_box}

\begin{claim_box}
\textbf{Consequence (a sequence of row operations):} If the
successive operations are represented by $E_1,\ldots,E_k$, then the
final matrix is
\[
E_k\cdots E_2E_1A.
\]
The first operation appears nearest to $A$, on the right of the other
elementary factors. This follows by applying the theorem at each step
and using associativity. The product $Q=E_k\cdots E_1$ is invertible,
with inverse $Q^{-1}=E_1^{-1}\cdots E_k^{-1}$, as direct cancellation
in both products verifies.

For an augmented matrix, the same operation acts on every column:
\[
E[A\mid b]=[EA\mid Eb].
\]
Since $E$ is invertible, $Ax=b$ if and only if $EAx=Eb$.
This gives a matrix explanation of why elementary row operations on
the entire augmented matrix preserve the solution set.
\end{claim_box}

\subsection{Homogeneous and nonhomogeneous linear systems}

\begin{definition_box}
\textbf{Definition:} A system of $m$ linear equations in $n$
unknowns over $\mathbb{F}$ can be written as
\[
Ax=b,
\qquad A\in\mathbb{F}^{m\times n},\quad
x\in\mathbb{F}^n,\quad b\in\mathbb{F}^m.
\]
Here $A=(a_{ij})$ is the given coefficient matrix, $b$ is the given
right-hand side, and $x$ is the unknown vector. The $i$th equation is
$\sum_{j=1}^n a_{ij}x_j=b_i$. The system is \emph{homogeneous} if
$b=0$, and \emph{nonhomogeneous} if $b\neq0$. A system is
\emph{consistent} if it has at least one solution, and
\emph{inconsistent} otherwise.
\end{definition_box}

\begin{claim_box}
\textbf{Definition (augmented matrix):} The \emph{augmented matrix}
of $Ax=b$ records the coefficients and right-hand side together:
\[
[A\mid b]=\left(\begin{array}{cccc|c}
a_{11}&a_{12}&\cdots&a_{1n}&b_1\\
\vdots&\vdots&\ddots&\vdots&\vdots\\
a_{m1}&a_{m2}&\cdots&a_{mn}&b_m
\end{array}\right).
\]
Each row represents one equation; the vertical line separates the
coefficient columns from the right-hand side.
\end{claim_box}

\begin{claim_box}
\textbf{Elementary row operations:} The following operations on
the \emph{entire} augmented matrix preserve the solution set:
\begin{enumerate}
    \item Interchange two rows: $R_i\leftrightarrow R_j$.
    \item Multiply a row by a nonzero scalar:
    $R_i\leftarrow cR_i$, where $c\neq0$.
    \item Add a scalar multiple of another row to a row:
    $R_i\leftarrow R_i+cR_j$, where $i\neq j$.
\end{enumerate}
\end{claim_box}

\begin{claim_box}
\textbf{Theorem (row operations preserve solutions):} Suppose
$[A'\mid b']$ is obtained from $[A\mid b]$ by a finite sequence of
elementary row operations. Then
\[
\{x\in\mathbb{F}^n:Ax=b\}
=\{x\in\mathbb{F}^n:A'x=b'\}.
\]

\textit{Proof:} It suffices to prove the claim for one operation.
Write $L_i(x)=\sum_{k=1}^n a_{ik}x_k$, so the original $i$th
equation is $L_i(x)=b_i$. Fix any candidate vector $x\in\mathbb{F}^n$.
\begin{enumerate}
    \item \emph{Interchanging rows:} Swapping rows $i$ and $j$ only
    changes the order of the equations. Thus $x$ satisfies all the
    new equations if and only if it satisfies all the original ones.
    \item \emph{Multiplying a row by a nonzero scalar:} The only
    changed equation is $cL_i(x)=cb_i$. Since $c\neq0$ has an inverse
    in $\mathbb{F}$,
    \[
    L_i(x)=b_i\quad\Longleftrightarrow\quad cL_i(x)=cb_i.
    \]
    All other equations are unchanged, so the solution set is unchanged.
    \item \emph{Adding a multiple of another row:} The operation
    $R_i\leftarrow R_i+cR_j$, with $i\neq j$, replaces equation $i$ by
    \[
    L_i(x)+cL_j(x)=b_i+cb_j
    \]
    and leaves equation $j$ and every other equation unchanged.
    If $x$ solves the original system, adding the two corresponding
    equalities shows that it satisfies the replacement equation.
    Conversely, if $x$ solves the new system, it satisfies both the
    replacement equation and $L_j(x)=b_j$. Subtracting $c$ times the
    latter gives
    \[
    L_i(x)=(b_i+cb_j)-cb_j=b_i.
    \]
    Hence $x$ also solves the original system.
\end{enumerate}
Each elementary operation therefore preserves exactly the same
solutions. Applying this conclusion at every step proves the result
for any finite sequence. \hfill$\square$
\end{claim_box}

\begin{remark_box}
\textbf{Caution:} Perform each operation on the entire augmented
row, including the right-hand side. The nonzero condition in row
scaling is essential: multiplying the equation $x=1$ by zero gives
$0=0$, which admits every $x$ and loses the original constraint.
\end{remark_box}

\begin{definition_box}
\textbf{Definition (pivot entries, positions, and columns):}
In a row echelon matrix, the leading entry of each nonzero row is
called a \emph{pivot entry}. Its location is a \emph{pivot position},
and a column containing one is a \emph{pivot column}. A pivot can be
any nonzero scalar in REF\@; in RREF every pivot is $1$ and all other
entries in its column are zero. Not every nonzero entry is a pivot.
For an arbitrary coefficient matrix $A$, its pivot columns mean the
column indices containing pivots after row reduction.
\end{definition_box}

\begin{claim_box}
\textbf{Pivot variables and free variables:} For a consistent
system, a variable $x_j$ is a \emph{pivot variable} if column $j$ of
the coefficient matrix has a pivot after row reduction. Otherwise it
is a \emph{free variable}, to which an arbitrary parameter can be
assigned. The final column of $[A\mid b]$ is not a variable column:
a pivot there instead signals an equation $0=d$ with $d\neq0$, hence
inconsistency.

For example, the RREF augmented matrix
\[
\left(\begin{array}{ccc|c}1&0&2&5\\0&1&-1&-1\end{array}\right)
\]
has pivot entries $1$ in positions $(1,1)$ and $(2,2)$. Its pivot
variables are $x_1,x_2$, and $x_3$ is free. The entries $2$ and $-1$
in column $3$ are not pivots. The equations are
$x_1=5-2x_3$ and $x_2=-1+x_3$.
\end{claim_box}

\begin{claim_box}
\textbf{Gaussian elimination:}
\begin{enumerate}
    \item Form $[A\mid b]$. Starting with the leftmost coefficient
    column, find a nonzero entry in or below the current row. Swap
    rows to place it in the current row, then use it to eliminate all
    entries below it. If no such entry exists, skip that column.
    Continue with the next row and the columns to the right until
    the matrix is in row echelon form.
    \item If a row has the form
    $\left(\begin{array}{ccc|c}0&\cdots&0&d\end{array}\right)$
    with $d\neq0$, it represents $0=d$. The system is inconsistent.
    A completely zero row imposes no condition.
    \item Otherwise, the variables whose coefficient columns contain
    pivots are \emph{pivot variables}; the remaining variables are
    \emph{free variables}. Assign arbitrary parameters in $\mathbb{F}$
    to the free variables and solve for the pivot variables from the
    bottom equation upward (\emph{back-substitution}).
    \item Alternatively, divide each pivot row by its pivot and
    eliminate entries above the pivots as well. This gives reduced
    row echelon form (\emph{Gauss--Jordan elimination}), from which
    the pivot variables can be read directly in terms of the free ones.
\end{enumerate}
\end{claim_box}

\begin{mainbox}{Theorem (pivot criteria for solutions)}
Let
$A\in\mathbb{F}^{m\times n}$, $b\in\mathbb{F}^m$, and let
$[R\mid d]$ be the RREF of $[A\mid b]$. Suppose $R$ has $r$ pivot
columns. Then:
\begin{enumerate}
    \item $Ax=b$ is consistent if and only if the augmented column
    contains no pivot, equivalently there is no row $[0\ \cdots\ 0\mid1]$.
    \item If the system is consistent, every assignment of the $n-r$
    free variables extends to exactly one solution.
    \item A consistent system has a unique solution if and only if
    $r=n$, that is, every variable column contains a pivot.
    \item The homogeneous system $Ax=0$ has a nonzero solution if and
    only if $r<n$.
\end{enumerate}
\end{mainbox}

\begin{claim_box}
\textit{Proof:} Row operations preserve solutions, so we may work with
$Rx=d$. A pivot in the augmented column gives the equation $0=1$,
making a solution impossible.

Conversely, suppose there is no such pivot. Every nonzero row then
has a pivot in a coefficient column. Let the pivot columns be
$p_1,\ldots,p_r$ in row order, and let $F$ be the set of the remaining
column indices. Since each pivot column has a single $1$ and zeros
elsewhere, row $i$ gives
\[
x_{p_i}=d_i-\sum_{j\in F}R_{ij}x_j,
\qquad 1\leq i\leq r.
\]
All remaining rows say $0=0$. Assign arbitrary values $x_j=t_j$ for
$j\in F$; these formulas determine the pivot variables uniquely and
satisfy every equation. This proves consistency and the second claim.

If $r=n$, there are no free variables, so the formulas give just one
solution. If $r<n$, assigning all free variables zero or assigning
one of them $1$ gives distinct solutions. Thus a consistent system
is unique exactly when $r=n$.

For $Ax=0$, apply the same sequence of row operations used on $[A\mid b]$.
The coefficient matrix again becomes $R$, while the zero right-hand
side remains zero throughout.
Thus the homogeneous system is always consistent, and if $r=n$ its
only solution is zero. If $r<n$, setting one free variable to $1$
produces a nonzero solution. \hfill$\square$
\end{claim_box}

\begin{claim_box}
\textbf{Proposition (pivot count equals rank):} The number $r$
of pivots in a row echelon form of $A$ equals $\operatorname{rank}A$.
Consequently, a consistent system has $n-\operatorname{rank}A$ free
variables.

\textit{Proof:} Further reducing an echelon form to RREF does not
change the number of pivots. Let $R$ be the resulting coefficient
matrix. Elementary row operations on vectors in $\mathbb{F}^m$ are
invertible linear maps: their inverses are the reverse row operations.
Their composition is an isomorphism $E:\mathbb{F}^m\longrightarrow
\mathbb{F}^m$, and
\[
Rx=E(Ax)\qquad\text{for all }x\in\mathbb{F}^n.
\]
Hence $\operatorname{range}T_R=E(\operatorname{range}T_A)$, so $R$
and $A$ have the same rank. The $r$ pivot columns of $R$ are the first
$r$ standard coordinate columns of $\mathbb{F}^m$. They are linearly
independent and span all columns of $R$, because every entry below
row $r$ is zero. Thus $\operatorname{rank}R=r$, proving the claim.
\hfill$\square$
\end{claim_box}

\begin{claim_box}
\textbf{Corollary:} We have
\[
Ax=b\text{ is consistent}
\quad\Longleftrightarrow\quad
\operatorname{rank}[A\mid b]=\operatorname{rank}A.
\]
Moreover, every $b\in\mathbb{F}^m$ gives a consistent system if and
only if $A$ has a pivot in every row after row reduction.

\textit{Proof:} In the reduced augmented matrix, the augmented column
adds a pivot exactly when the system is inconsistent. The rank formula
therefore follows from the pivot-count proposition. Finally, every
right-hand side is attainable exactly when $\operatorname{range}T_A$
is all of $\mathbb{F}^m$, equivalently $\operatorname{rank}A=m$.
By the proposition, this means there are $m$ pivots, one in every row.
\hfill$\square$
\end{claim_box}

\begin{remark_box}
\textbf{Note:} For a consistent system,
the free-variable assignments give a bijection between the solution
set and $\mathbb{F}^{n-r}$. Thus if $r<n$ there are infinitely many
solutions over an infinite field; over a finite field with $q$ elements
there are exactly $q^{n-r}$ solutions.
\end{remark_box}

\begin{example_box}
\textbf{Example (solving by row reduction):} Consider
\[
x+2y=3,\qquad -y+z=1.
\]
The augmented matrix reduces as follows:
\[
\begin{aligned}
\left(\begin{array}{ccc|c}1&2&0&3\\0&-1&1&1\end{array}\right)
&\xrightarrow{R_2\leftarrow -R_2}
\left(\begin{array}{ccc|c}1&2&0&3\\0&1&-1&-1\end{array}\right)\\[4pt]
&\xrightarrow{R_1\leftarrow R_1-2R_2}
\left(\begin{array}{ccc|c}1&0&2&5\\0&1&-1&-1\end{array}\right).
\end{aligned}
\]
The pivot variables are $x,y$, and $z$ is free. Setting $z=t$ gives
\[
x=5-2t,\qquad y=t-1,\qquad z=t,
\]
so all solutions are
$(x,y,z)=(5,-1,0)+t(-2,1,1)$ for $t\in\mathbb{R}$.
With $s=t-1$, this is equivalently
$(3,0,1)+s(-2,1,1)$, the form used below.
\end{example_box}

\begin{example_box}
\textbf{Example (detecting inconsistency):} For
$x+y=1$ and $2x+2y=3$,
\[
\left(\begin{array}{cc|c}1&1&1\\2&2&3\end{array}\right)
\xrightarrow{R_2\leftarrow R_2-2R_1}
\left(\begin{array}{cc|c}1&1&1\\0&0&1\end{array}\right).
\]
The last row says $0=1$, so there is no solution.
\end{example_box}

\begin{claim_box}
\textbf{Homogeneous systems:} The system $Ax=0$ always has the
\emph{trivial solution} $x=0$. Its complete solution set is
\[
\{x\in\mathbb{F}^n:Ax=0\}=\ker T_A,
\qquad T_A(x)=Ax,
\]
so it is a subspace of $\mathbb{F}^n$. In particular, every linear
combination of homogeneous solutions is another homogeneous solution.
Writing $r=\operatorname{rank}A=\dim(\operatorname{range}T_A)$,
rank-nullity gives
\[
\dim(\ker T_A)=n-r.
\]
Thus a nontrivial solution exists exactly when $r<n$.
\end{claim_box}

\begin{claim_box}
\textbf{Nonhomogeneous systems and consistency:} The system
$Ax=b$ has a solution exactly when
\[
b\in\operatorname{range}T_A.
\]
Equivalently, $b$ must be a linear combination of the columns of $A$,
because $Ax=x_1 A_1+\cdots+x_n A_n$, where $A_j$ is the $j$th column.
Unlike a homogeneous system, a nonhomogeneous system need not be
consistent.
\end{claim_box}

\begin{mainbox}{Proposition (the complete solution set)}
Suppose $x_p$
is one particular solution of $Ax=b$. Then all solutions are given by
\[
x=x_p+h,\qquad h\in\ker T_A.
\]
In other words, the solution set is $x_p+\ker T_A$.
\end{mainbox}

\begin{claim_box}
\textit{Proof:} Since $Ax_p=b$, for any $x\in\mathbb{F}^n$,
\[
Ax=b\quad\Longleftrightarrow\quad A(x-x_p)=0
\quad\Longleftrightarrow\quad x-x_p\in\ker T_A.
\]
This proves both that every solution has the stated form and that every
vector of that form is a solution. \hfill$\square$
\end{claim_box}

\begin{claim_box}
\textbf{Interpretation:} A consistent solution set is a translate
of the null space, also called an \emph{affine subspace}. If $b\neq0$,
it is not a vector subspace: it does not contain $0$, since $A0=0\neq b$.
The difference of any two solutions lies in the null space, but the sum
of two solutions need not be a solution of the same system.
\end{claim_box}

\begin{example_box}
\textbf{Example (a family of solutions):} Over $\mathbb{R}$,
consider the homogeneous system
\[
x+2y=0,\qquad -y+z=0.
\]
Its solutions are $(x,y,z)=t(-2,1,1)$, where $t\in\mathbb{R}$.
Now change the right-hand side:
\[
x+2y=3,\qquad -y+z=1.
\]
A particular solution is $(3,0,1)$. Thus all solutions are
\[
(x,y,z)=(3,0,1)+t(-2,1,1),\qquad t\in\mathbb{R}.
\]
The homogeneous solution set is a line through the origin; the
nonhomogeneous solution set is a parallel line that does not pass
through the origin. Both have one free parameter.
\end{example_box}

\begin{example_box}
\textbf{Examples (no solution or a unique solution):} Over
$\mathbb{R}$, the system
\[
x+y=1,\qquad 2x+2y=3
\]
is inconsistent: doubling the first equation would give $2x+2y=2$,
contradicting the second. In contrast,
\[
x+y=3,\qquad x-y=1
\]
has the unique solution $(x,y)=(2,1)$. Its associated homogeneous system
has only the zero solution.
\end{example_box}

\begin{claim_box}
\textbf{Consequences of rank-nullity:}
\begin{enumerate}
    \item A consistent system has a unique solution if and only if
    $\ker T_A=\{0\}$, equivalently $r=n$. Otherwise its solutions have
    $n-r$ free parameters: if $h_1,\ldots,h_{n-r}$ is a basis of the
    kernel, every solution is uniquely expressible as
    \[
    x=x_p+c_1 h_1+\cdots+c_{n-r}h_{n-r},\qquad c_i\in\mathbb{F}.
    \]
    Over an infinite field such as $\mathbb{R}$ or $\mathbb{C}$, a
    consistent system with $r<n$ therefore has infinitely many solutions.
    \item If $m<n$, then $r\leq m<n$, so the homogeneous system has
    nonzero solutions. A nonhomogeneous system may be inconsistent,
    but if it is consistent, its solution cannot be unique.
    \item Every right-hand side $b\in\mathbb{F}^m$ gives a consistent
    system exactly when $T_A$ is surjective, equivalently $r=m$.
    In particular, if $m>n$, some right-hand sides have no solution;
    this does not mean that every right-hand side is inconsistent.
\end{enumerate}
\end{claim_box}

\end{document}